% !TeX root = ../../Book1.tex
\chapter{Lipschitz 映射、近似微分与绝对连续函数}
\label{b1:ch:16}
\BookChapterNumber{b1:ch:16}
% 本章摘要（不计入编号）
\begin{chapterabstract}
Lipschitz 条件只比较两点的距离，却能迫使欧氏空间中的映射几乎处处具有线性近似。
本章证明 Rademacher 定理，说明从直线上的导数走向全微分时，
为什么还需要平移测试和方向上的一致控制。
随后用近似微分处理可测映射的局部结构，并给出几乎处处近似可微映射的可数 Lipschitz 分片。
实值绝对连续函数的微积分基本定理为后续 Sobolev 理论提供一维基础；
选读节再把速度与长度的讨论推广到一般度量空间，分清向量导数与度量导数。
\end{chapterabstract}

% 本章目标（不计入编号）
\begin{learninggoals}
\begin{itemize}
\item\label{b1:item:ch16-goal-extension}掌握 Lipschitz 常数、双边距离估计及实值延拓，理解像集测度估计的适用范围。
\item\label{b1:item:ch16-goal-rademacher}能说明 Rademacher 证明中一维积分表示、方向导数线性关系及有限方向网各自解决的问题。
\item\label{b1:item:ch16-goal-approximate}区分普通微分与近似微分，理解可数 Lipschitz 分片的结论及其不能控制的部分。
\item\label{b1:item:ch16-goal-ac}掌握实值绝对连续函数、全变差与微积分基本定理，理解可积导数怎样恢复函数的全部增量。
\item\label{b1:item:ch16-goal-curves}在选读部分理解绝对连续曲线的可积速度控制、度量导数与长度公式，并辨认有限维结论的边界。
\end{itemize}
\end{learninggoals}

\secref{b1:sec:1602}是本章主线：\chapref{b1:ch:15}由覆盖构造的 Hausdorff 测度及其 Lipschitz 映射估计尚不足以精确计算一个映射如何改变体积，
Rademacher 定理提供的几乎处处微分矩阵才使这种计算有了局部对象。
近似微分一节可以在首次读完该定理后继续阅读；
\secref{b1:sec:1604}则回答一维函数何时能由导数积分恢复；
\secref{b1:sec:1605}标为选读，把这些结果推广到一般度量值曲线。

\ProseParagraphBreak

Simon 的《Lectures on Geometric Measure Theory》%
\cite[\S5]{Simon1983Lectures}提供本章可微性证明的主要参照，
Fremlin 的《Measure Theory》%
\cite[\S\S225 and 262]{Fremlin2016Measure}可用于对照绝对连续性与 Lipschitz 映射的测度论处理。
本书根据前文已得的结果安排两项证明：一维 Lipschitz 积分表示先由 \(L^1\) 对偶得到，
绝对连续函数的积分表示则从变差与 Radon--Nikodym 定理建立。

% !TeX root = ../../Book1.tex
\section{Lipschitz 映射}
\label{b1:sec:1601}

上一章只用距离控制就得到了 Hausdorff 测度的映射估计。
本章进一步说明，这种不涉及导数的距离条件，为什么仍能产生几乎处处的线性近似。
证明 Rademacher 定理以前，先介绍 Lipschitz 常数、延拓和参数化的基本性质。
其中实值延拓尤其有用：它允许把定义在任意子集上的函数变成全空间函数，
而不要求原来的定义域有光滑边界。

\ProseParagraphBreak

% 实际读取：Fremlin2016Measure 作者官方chap26.tex，262A、262B(a)(b)完整正文。
% 本节把标量延拓写成上下两种包络，并单列向量逐坐标延拓的sqrt(m)常数。
% Hausdorff估计引用本书15.1已证结果；像的可测性不借尚未建立的描述集合论。
关于这些基本操作，可对读 Fremlin 的《Measure Theory》%
\cite[\S262]{Fremlin2016Measure}。本节把延拓所需的有限值检查与向量常数写出，
后面的可微性证明将直接使用它们。

\subsection{Lipschitz 常数}
\label{b1:subsec:160101}

\begin{definition}\label{b1:def:lipschitz-map}
设 \((X,\mathrm{d}_X)\)、\((Y,\mathrm{d}_Y)\) 是度量空间，\(A\subseteq X\)。
若存在 \(L<\infty\)，使
\[
 \mathrm{d}_Y\bigl(f(x),f(y)\bigr)\leq L\cdot \mathrm{d}_X(x,y)
 \quad(x,y\in A),
\]
则称 \(f:A\rightarrow Y\) 为 \term[lipschitz-yingshe]{Lipschitz 映射}[Lipschitz map]，
并称每个满足此式的非负数 \(L\) 为它的一个 Lipschitz 常数。
记最小常数为
\[
 \operatorname{Lip}(f;A)
 \coloneqq\sup_{\substack{x,y\in A\\x\ne y}}
   \frac{\mathrm{d}_Y\bigl(f(x),f(y)\bigr)}{\mathrm{d}_X(x,y)};
\]
当 \(A\) 至多有一个点时，规定此上确界为零。
\end{definition}
\symidx{\(\operatorname{Lip}(f;A)\)：最小 Lipschitz 常数}

最小常数确实仍满足距离不等式，而不是只给出可取常数的下确界。
Lipschitz 映射一致连续，并把 Cauchy 列映成 Cauchy 列。
若目标空间完备，则它能唯一连续延拓到 \(\overline A\)，且 Lipschitz 常数不变：
对 \(x\in\overline A\) 取 \(A\) 中收敛到 \(x\) 的序列，
以像序列的极限定义延拓值；交错两条逼近列保证定义无关，再对不等式取极限。
这一结论只是在闭包上延拓，不能与稍后的全空间延拓混为一谈。

\ProseParagraphBreak

若 \(f\)、\(g\) 可以复合，则
\[
 \operatorname{Lip}(g\circ f;A)
 \leq\operatorname{Lip}(g;f(A))\operatorname{Lip}(f;A).
\]
对向量值函数，和与标量乘法的常数估计由三角不等式直接得到。
特别地，点到非空集合 \(E\) 的距离函数满足
\[
 \bigl|\operatorname{dist}(x,E)-\operatorname{dist}(y,E)\bigr|
 \leq \mathrm{d}_X(x,y),
\]
因为先对 \(z\in E\) 使用三角不等式，再取下确界并交换 \(x,y\)。
距离函数即使没有普通导数，也始终具有这个定量控制。

\ProseParagraphBreak

若对定义域内每一点，都存在一个邻域，使函数在该邻域内 Lipschitz，
则称其局部 Lipschitz。局部条件不自动给出一个对整个定义域有效的常数；
例如 \(x\mapsto x^2\) 在 \(\mathbb R\) 上局部 Lipschitz，却不是全局 Lipschitz。
对于开集上的局部 Lipschitz 映射，可取可数个小球覆盖定义域，
使函数在每个球上都有一个常数。这个可数化保证后面逐球取得的零测例外仍可合并为零测集。

\begin{proposition}\label{b1:prop:bounded-derivative-lipschitz}
设 \(\Omega\subseteq\mathbb R^d\) 是开凸集，
\(f:\Omega\rightarrow\mathbb R^m\) 在每一点可微，且
\(\|Df(x)\|\leq L\) 对所有 \(x\in\Omega\) 成立。
则 \(f\) 为 \(L\)-Lipschitz 映射。
\end{proposition}
\begin{proof}
任取 \(x,y\in\Omega\) 及 \(\mathbb R^m\) 中的单位向量 \(v\)。
对实函数 \(t\mapsto\langle v,f(x+t(y-x))\rangle\) 使用一维中值定理，
其导数的绝对值不超过 \(L|y-x|\)，故
\[
 |\langle v,f(y)-f(x)\rangle|\leq L|y-x|.
\]
对单位向量 \(v\) 取上确界即得结论。凸性保证连接两点的整条线段都在定义域内。
\end{proof}

若定义域不连通，“导数处处为零”也可能只说明函数在各个连通部分上常值。
例如在 \((-1,0)\cup(0,1)\) 上，左右两部分分别取值 \(0\)、\(1\)，
函数处处可微且导数为零，却不是 Lipschitz 的。
这说明从局部导数信息恢复全局距离估计，需要检查定义域中的连接条件。

\subsection{双 Lipschitz 映射}
\label{b1:subsec:160102}

\secref{b1:sec:1502}已经定义双 Lipschitz 映射。
它除距离上界外还具有严格正的下界
\[
 a|x-y|\leq|f(x)-f(y)|\leq b|x-y|,
 \quad0<a\leq b<\infty.
\]
左侧立即保证单射，也保证像上的逆映射具有常数 \(a^{-1}\)。
因此这是一种同时控制两边距离的参数化，而不是一般的连续一一对应。

\ProseParagraphBreak

线性情形最能说明这个差别。若 \(T:\mathbb R^k\rightarrow\mathbb R^d\) 单射，
则连续函数 \(v\mapsto|Tv|\) 在单位球面上有正的最小值 \(\sigma\)，从而
\[
 \sigma|x-y|\leq|T(x-y)|\leq\|T\|\,|x-y|.
\]
若 \(T\) 有非零核向量，则任何正的下界都不可能成立。
非线性映射在某一点有单射微分，并不单凭这一点就保证它在某邻域内双 Lipschitz；
微分在附近是否受控同样重要。

\ProseParagraphBreak

一个经常使用的充分条件是对线性单射作小 Lipschitz 扰动。
若 \(|Tv|\geq\sigma|v|\)，且 \(g:A\rightarrow\mathbb R^d\) 满足
\(\operatorname{Lip}(g;A)=\eta<\sigma\)，则对 \(f(x)=Tx+g(x)\) 有
\[
 (\sigma-\eta)|x-y|
 \leq|f(x)-f(y)|
 \leq(\|T\|+\eta)|x-y|.
\]
两侧分别来自反向三角不等式和三角不等式。
这里的小量是 \(g\) 的距离差商，而不仅是它的函数值。
即使函数值在一个小邻域内很小，其振荡仍可能破坏单射性。

\subsection{Lipschitz 像}
\label{b1:subsec:160103}

设 \(f:A\subseteq\mathbb R^k\rightarrow\mathbb R^d\) 为 Lipschitz 映射。
若 \(A\) 紧，则 \(f(A)\) 紧；若 \(A\) 是可数个紧集的并，则像也是可数个紧集的并。
但对一般 Borel 定义域，不能仅凭连续性就断言其像是 Borel 集。
本书在尚未证明像的可测性时，对 \(\mathcal H^s\bigl(f(A)\bigr)\) 一律按外测度解释。

\ProseParagraphBreak

参数化往往先在一个不规则的子集上给出。
以下实值延拓把这类定义域问题与可微性问题分开处理。

\begin{theorem}[McShane 延拓]\label{b1:thm:mcshane-extension}
设 \((X,\mathrm{d})\) 为度量空间，\(A\subseteq X\) 非空，
\(f:A\rightarrow\mathbb R\) 为 \(L\)-Lipschitz 函数。
对每个 \(x\in X\)，定义
\[
 F_-(x)=\sup_{a\in A}\bigl(f(a)-L\cdot \mathrm{d}(x,a)\bigr),
 \quad
 F_+(x)=\inf_{a\in A}\bigl(f(a)+L\cdot \mathrm{d}(x,a)\bigr).
\]
则这两个函数处处有限，均为 \(f\) 的 \(L\)-Lipschitz 延拓，并满足 \(F_-\leq F_+\)。
\end{theorem}
\term[mcshane-yantuo]{McShane 延拓定理}[McShane extension theorem]的一个特点是，
原定义域既不必闭，也不必有界；因此证明中必须先检查上、下确界不会为无穷。

\begin{proof}
固定 \(a_0\in A\)。由 \(f\) 的 Lipschitz 条件和三角不等式，
对所有 \(a\in A\) 都有
\[
 f(a)+L\cdot \mathrm{d}(x,a)
 \geq f(a_0)-L\cdot \mathrm{d}(a,a_0)+L\cdot \mathrm{d}(x,a)
 \geq f(a_0)-L\cdot \mathrm{d}(x,a_0).
\]
而取 \(a=a_0\) 又给出有限上界，故 \(F_+(x)\) 有限。
同理，\(F_-(x)\) 也有限。
若 \(x\in A\)，Lipschitz 条件给出 \(F_+(x)\geq f(x)\)，
而候选点 \(a=x\) 给出反向不等式；对 \(F_-\) 一样，故两者都延拓 \(f\)。

对每个 \(a\in A\)，有
\[
 f(a)+L\cdot \mathrm{d}(x,a)\leq f(a)+L\cdot \mathrm{d}(y,a)+L\cdot \mathrm{d}(x,y).
\]
取下确界得 \(F_+(x)\leq F_+(y)+L\cdot \mathrm{d}(x,y)\)，交换两点得到 Lipschitz 条件。
对 \(F_-\) 取上确界可作同样证明。
最后，对任意 \(a,b\in A\)，
\[
 f(a)-L\cdot \mathrm{d}(x,a)
 \leq f(b)+L\cdot \mathrm{d}(a,b)-L\cdot \mathrm{d}(x,a)
 \leq f(b)+L\cdot \mathrm{d}(x,b).
\]
先对 \(a\) 取上确界，再对 \(b\) 取下确界，即得 \(F_-\leq F_+\)。
\end{proof}

对 \(f:A\subseteq\mathbb R^k\rightarrow\mathbb R^d\)，逐坐标使用该定理，
即可得到全空间的向量值延拓 \(\widetilde f\)，并有
\[
 \operatorname{Lip}(\widetilde f;\mathbb R^k)\leq\sqrt d\,\operatorname{Lip}(f;A).
\]
因为每个坐标的差至多为 \(L|x-y|\)，把 \(d\) 个坐标的平方相加就得到这个常数。
这里没有证明向量延拓保持最优常数；后面的局部可微性论证只需要一个有限常数，
故逐坐标构造已经足够。

\ProseParagraphBreak

作为典型的参数化，若 \(u:A\subseteq\mathbb R^k\rightarrow\mathbb R^m\)
是 \(L\)-Lipschitz 的，则它的图像映射 \(G(x)=(x,u(x))\) 满足
\[
 |x-y|\leq|G(x)-G(y)|\leq\sqrt{1+L^2}\,|x-y|.
\]
图像因而与定义域双 Lipschitz 等价，逆映射就是坐标投影。
这给出了不要求函数光滑的曲面片模型，后面将用其微分描述局部切向信息。

\subsection{Hausdorff 测度估计}
\label{b1:subsec:160104}

\cref{b1:prop:hausdorff-lipschitz}已给出核心估计：
对 \(s>0\)、任意 \(E\subseteq A\)，
\[
 \mathcal H^s\bigl(f(E)\bigr)\leq L^s\mathcal H^s(E).
\]
本节不重复覆盖证明，而列出它与本章问题直接相关的后果。
若 \(f\) 为双 Lipschitz 映射，还同时有
\[
 a^s\mathcal H^s(E)\leq\mathcal H^s\bigl(f(E)\bigr)
 \leq b^s\mathcal H^s(E).
\]
因此零测性和有限正测度性质在这种参数化下保持。

\ProseParagraphBreak

特别地，Lipschitz 映射 \(f:A\subseteq\mathbb R^d\rightarrow\mathbb R^d\)
把 Lebesgue 零测集映成 Lebesgue 零测集。
这使用了同维的 \(\mathcal H^d=m_d^*\)。
若目标是更高维空间，最直接的结论仍是像的 \(d\) 维 Hausdorff 测度为零，
不能把下标随意改成目标空间维数来代替这一更强信息。
对局部 Lipschitz 映射，以可数个具有有限常数的局部块覆盖定义域，再取可数并，同样成立。

\ProseParagraphBreak

图像模型给出一个具体的双边比较：
\[
 \mathcal H^k(E)\leq
 \mathcal H^k\bigl(\{(x,u(x)):x\in E\}\bigr)
 \leq(1+L^2)^{k/2}\mathcal H^k(E).
\]
但这还不是面积公式。它用一个全局常数控制所有点，
没有记录不同位置的倾斜程度，也没有处理一般映射的多重覆盖。
\chapref{b1:ch:17}将以微分矩阵给出更精确的计算。

\ProseParagraphBreak

对一条 \(L\)-Lipschitz 曲线 \(\gamma:[a,b]\rightarrow\mathbb R^d\)，
已经有
\[
 \mathcal H^1\bigl(\gamma([a,b])\bigr)\leq L(b-a).
\]
而对任意连续曲线，像的 \(\mathcal H^1\) 测度至少为其直径：
取像中的 \(p,q\)，距离函数 \(z\mapsto|z-p|\) 是 \(1\)-Lipschitz 的，
由连续性，其在曲线像上的取值包含区间 \([0,|q-p|]\)；
使用测度映射估计及 \(\mathcal H^1([0,|q-p|])=|q-p|\)，
再对 \(p,q\) 取上确界即可。
这些估计只涉及像作为一个集合的大小。
若一条曲线反复经过同一线段，像并不会因此变大，但沿参数走过的总长度会增加。
\secref{b1:sec:1605}将把两种量分清，并用度量导数计算后者。

% !TeX root = ../../Book1.tex
\section{Rademacher 定理}
\label{b1:sec:1602}

Lipschitz 条件控制任意两点之间的增量，却没有直接给出增量的一阶主部。本节证明：在欧氏空间中，这个一阶主部几乎处处存在。论证采用 Leon Simon 的《Lectures on Geometric Measure Theory》\cite[5.2 THEOREM, pp. 30--32]{Simon1983Lectures}中的方向导数路线：先沿直线取得导数，再借平移下的积分恒等式识别它们之间的线性关系，最后把逐方向的极限提升为全微分。

% 来源：实际读取 Simon1983Lectures §5.2，印刷 pp.30--32；
% 对照 Fremlin2016Measure 官方 chap26.tex 的 262Q。
% Simon 在一维处调用绝对连续理论；本节改用已证 L1 对偶表示补齐此桥，
% 不预借16.4的一般绝对连续函数定理。紧支集测试分离用14.4的近似恒等核。

\subsection{一维 Lipschitz 函数}
\label{b1:subsec:160201}

先考察闭区间上的函数。此时不只是差商几乎处处收敛：函数的全部增量还可以由一个有界密度的积分恢复。这一点将保证沿直线取得的导数保留原函数的信息。

\begin{theorem}\label{b1:thm:one-dimensional-lipschitz-ftc}
设 \(-\infty<a<b<\infty\)，\(m\ge1\) 为整数，\(0\le L<\infty\)，且 \(f:[a,b]\rightarrow\mathbb R^m\) 为 \(L\)-Lipschitz 映射。则存在可测函数 \(g:(a,b)\rightarrow\mathbb R^m\)，使 \(|g|\leq L\) 几乎处处成立，并且对全部 \(a\leq s<t\leq b\)，
\[
 f(t)-f(s)=\int_s^t g(u)\,\mathrm du.
\]
这里的向量积分按分量定义。函数 \(f\) 在 \((a,b)\) 的几乎每个点可导，且 \(f'=g\) 几乎处处；上述积分密度在几乎处处相等的意义下唯一。
\end{theorem}
\begin{proof}
先设 \(m=1\)。对区间阶梯函数
\[
 \phi=\sum_{i=1}^N c_i\cdot\mathbf1_{(t_{i-1},t_i]},
 \quad a=t_0<t_1<\cdots<t_N=b,
\]
定义
\[
 \Lambda(\phi)=\sum_{i=1}^N c_i\bigl(f(t_i)-f(t_{i-1})\bigr).
\]
细分区间不改变右端，因为相邻增量相加会消去中间端点值。两种表示取共同细分后，在每个正长度开区间上的系数必须一致；有限个端点不影响 \(L^1(a,b)\) 的等价类。因此 \(\Lambda\) 是这个阶梯函数子空间上的良定义线性泛函，而且
\[
 |\Lambda(\phi)|
 \leq L\sum_{i=1}^N|c_i|(t_i-t_{i-1})
 =L\|\phi\|_1.
\]
这类阶梯函数在 \(L^1(a,b)\) 中稠密。事实上，把任意 \(L^1(a,b)\) 函数在区间外延为零，由\cref{b1:prop:lp-continuous-dense}先用 \(\mathbb R\) 上的连续紧支集函数作 \(L^1\) 逼近；再利用其在 \([a,b]\) 上的一致连续性，用有限区间阶梯函数一致逼近其限制。

由\cref{b1:prop:continuous-extension-closure}，\(\Lambda\) 唯一延拓为整个 \(L^1(a,b)\) 上的连续线性泛函，范数不超过 \(L\)。再用\cref{b1:thm:l1-dual-sigma-finite}，得到实函数 \(g\in L^\infty(a,b)\)，满足
\[
 \|g\|_\infty\leq L,\quad
 \Lambda(\phi)=\int_a^b\phi(u)\cdot g(u)\,\mathrm du.
\]
测试 \(\mathbf1_{(s,t]}\) 就得到所需增量公式。这是对任意端点成立的等式，不是仅在删去一组端点后成立的表示。

取 \(g\) 的一个有限值代表。若 \(x\in(a,b)\) 是它的 Lebesgue 点，则对充分小的非零 \(h\)，无论 \(h\) 正负，都有
\[
 \left|\frac{f(x+h)-f(x)}h-g(x)\right|
 \leq\frac1{|h|}\int_{x-|h|}^{x+|h|}
                 |g(u)-g(x)|\,\mathrm du
 \longrightarrow0.
\]
右端是中心区间平均绝对偏差的两倍，故由\cref{b1:thm:lebesgue-point}趋零。这证明 \(f'=g\) 几乎处处，也说明任何给出同样增量公式的可积密度都必须几乎处处等于 \(f'\)。

对一般 \(m\)，逐个实分量应用已证结论，并删去有限个例外零测集的并。在剩余点，各分量差商同时收敛，因而向量差商收敛到 \(g\)。原来的 Lipschitz 估计直接给出
\[
 |g(x)|=\lim_{h\to0}\frac{|f(x+h)-f(x)|}{|h|}\leq L.
\]
于是得到欧氏范数下的界 \(L\)，而不只是从分量界拼出的 \(\sqrt mL\)。积分等式及唯一性也逐分量成立。
\end{proof}

例如 \(f(t)=|t|\) 的密度可取为 \(\operatorname{sgn}t\)，并在原点任意赋值。这个函数在原点有折角，但折角不妨碍积分恢复任意两个端点之间的增量。定理并不声称导数处处存在，更不声称导数连续。

\subsection{高维可微性}
\label{b1:subsec:160202}

在高维情形，需要先分清沿一条直线的导数与同时控制全部增量的微分。

\begin{definition}\label{b1:def:frechet-differentiability}
设 \(U\subseteq\mathbb R^n\) 为开集，\(f:U\rightarrow\mathbb R^m\)，且 \(x\in U\)。若存在一个线性映射 \(A:\mathbb R^n\rightarrow\mathbb R^m\)，使
\[
 \lim_{\substack{h\to0\\h\ne0}}
       \frac{|f(x+h)-f(x)-Ah|}{|h|}=0,
\]
则称 \(f\) 在 \(x\) 处\term[kewei]{可微}[differentiable]，称 \(A\) 为它在该点的\term[quan-weifen]{全微分}[total differential]，记作 \(Df(x)\)。若固定 \(v\in\mathbb R^n\)，且有限极限
\[
 D_vf(x)\coloneqq\lim_{t\to0}\frac{f(x+tv)-f(x)}t
\]
存在，则称其为 \(f\) 在 \(x\) 处沿 \(v\) 的\term[fangxiang-daoshu]{方向导数}[directional derivative]。
\end{definition}
\symidx{\(Df(x)\)：映射在一点的全微分}

全微分唯一：若 \(A,B\) 都满足定义，沿 \(h=tv\) 取极限便得到 \(Av=Bv\)，这对每个 \(v\) 成立。可微性还给出 \(D_vf(x)=Df(x)v\)。若 \(f\) 在 \(x\) 附近为 \(L\)-Lipschitz，则沿同样的增量取极限，得到
\[
 |Df(x)v|\leq L|v|,\quad \|Df(x)\|\leq L,
\]
其中 \(\|\cdot\|\) 是线性映射的算子范数。

\ProseParagraphBreak

反过来，所有方向导数都存在，并不足以给出一个全微分。考虑
\[
 F(x,y)=
 \begin{cases}
 \dfrac{xy^2}{x^2+y^2},&(x,y)\ne(0,0),\\[4pt]
 0,&(x,y)=(0,0).
 \end{cases}
\]
这个函数实际上是全局 Lipschitz 的。原点外的偏导数为
\[
 \partial_xF=\frac{y^2(y^2-x^2)}{(x^2+y^2)^2},
 \quad
 \partial_yF=\frac{2x^3y}{(x^2+y^2)^2},
\]
二者的绝对值都不超过 \(1\)，故梯度范数不超过 \(\sqrt2\)。不经过原点的线段可直接应用一元中值估计；若线段经过原点，则分成两段，并利用 \(F\) 在原点的连续性取极限。因此 \(F\) 的 Lipschitz 常数不超过 \(\sqrt2\)。然而，对非零方向 \(v=(v_1,v_2)\)，
\[
 D_vF(0,0)=\frac{v_1v_2^2}{v_1^2+v_2^2}.
\]
两个坐标方向的导数均为零，而方向 \((1,1)\) 的导数为 \(1/2\)，所以这些数不可能来自同一个线性映射。

\ProseParagraphBreak

Rademacher 定理的证明必须排除上述非线性现象在正测度集合上发生。为此，先给出一个以积分辨认局部可积函数的工具。以下 \(C_c^1(U)\) 表示具有紧支集、且一阶偏导数连续的实值函数。

\begin{lemma}\label{b1:lem:compact-tests-separate}
设 \(U\subseteq\mathbb R^n\) 是开集，积分针对 Lebesgue 测度。设 \(w\in L^1_{\mathrm{loc}}(U)\) 为实值函数。如果对每个 \(\phi\in C_c^1(U)\) 都有
\[
 \int_U w(y)\cdot\phi(y)\,\mathrm dy=0,
\]
那么 \(w=0\) 几乎处处。
\end{lemma}
\begin{proof}
取
\[
 \eta(z)=c(1-|z|^2)_+^2,
 \quad c^{-1}=\int_{B(0,1)}(1-|z|^2)^2\,\mathrm dz.
\]
归一化常数正且有限，\(\eta\) 非负、径向非增、积分为 \(1\)。由于球内的函数及其一阶偏导数在球面上都与外侧的零函数相接，\(\eta\in C_c^1(\mathbb R^n)\)。固定 \(x\in U\)，并取 \(\overline B(x,r)\subseteq U\)。函数
\[
 \phi_{x,r}(y)=r^{-n}\cdot \eta\left(\frac{y-x}r\right)
\]
属于允许的测试类，故其与 \(w\) 的积分为零。若 \(x\) 是 \(w\) 的 Lebesgue 点，\cref{b1:prop:approx-identity-lebesgue-point}的紧支集版本说明该积分在 \(r\downarrow0\) 时趋于 \(w(x)\)。几乎每个点都是这样的点，所以 \(w=0\) 几乎处处。
\end{proof}

\subsection{Rademacher 定理}
\label{b1:subsec:160203}

\term[rademacher-dingli]{Rademacher 定理}[Rademacher theorem]中的“几乎处处”指 \(n\) 维 Lebesgue 测度。它适用于任意有限维值域，不要求定义域中的函数已经具有连续偏导数。

\ProseParagraphBreak

证明分四步进行。先对每个固定方向沿平行直线应用一维结果，得到该方向上的几乎处处导数；
再用平移换元后的测试函数恒等式，识别不同方向导数之间的线性关系；
随后只保留可数稠密方向，从而把各方向的例外集统一成一个零测集；
最后用有限方向网与 Lipschitz 估计，把稠密方向上的差商收敛升级为任意小增量下的 Fr\'echet 微分。
也就是说，证明的主链是
\[
 \begin{gathered}
 \text{固定方向可导}
 \longrightarrow \text{平移测试识别线性关系}\\
 \Downarrow\\
 \text{可数稠密方向统一零集}
 \longrightarrow \text{有限方向网得到全微分}.
 \end{gathered}
\]

\cref{b1:fig:rademacher-local-picture}提醒我们，Lipschitz 图像可以有尖角；定理断言的不是处处可微，而是这类坏点只能落在零测例外集中。
\begin{figure}[htbp]
  \centering
  \begin{tikzpicture}[scale=1.05,>=stealth]
    \begin{scope}
      \draw[->] (-1.65,0)--(1.7,0);
      \draw[->] (0,-0.2)--(0,1.35);
      \draw[thick,domain=-1.45:1.45,samples=60] plot (\x,{0.45*\x*\x});
      \draw[dashed,domain=0.15:1.48,samples=2] plot (\x,{0.72*\x-0.288});
      \fill (0.8,0.288) circle (1.4pt) node[above left] {$x$};
      \node[below=0.35cm] at (0,-0.2) {线性主部存在};
    \end{scope}
    \begin{scope}[shift={(4.2,0)}]
      \draw[->] (-1.65,0)--(1.7,0);
      \draw[->] (0,-0.2)--(0,1.35);
      \draw[thick] (-1.4,1.05)--(0,0)--(1.4,1.05);
      \fill (0,0) circle (1.4pt) node[above right] {$x$};
      \node[below=0.35cm] at (0,-0.2) {尖角处无全微分};
    \end{scope}
  \end{tikzpicture}
  \caption[可微点与尖角]{可微点与尖角。左图的虚线表示一阶线性近似；右图说明 Lipschitz 条件仍允许出现不可微点。}
  \label{b1:fig:rademacher-local-picture}
\end{figure}

\begin{theorem}[Rademacher]\label{b1:thm:rademacher}
设 \(U\subseteq\mathbb R^n\) 为开集，且 \(f:U\rightarrow\mathbb R^m\) 局部 Lipschitz。则 \(f\) 在 \(U\) 的几乎每个点可微。若 \(f\) 在某个邻域内为 \(L\)-Lipschitz，则该邻域内每个可微点均满足 \(\|Df(x)\|\leq L\)。
\end{theorem}
\begin{proof}
先证明全空间上的实值 \(L\)-Lipschitz 函数。对每个固定的单位向量 \(v\)，记
\[
 \Delta_{v,t}f(x)=\frac{f(x+tv)-f(x)}t\quad(t\ne0).
\]
每个差商都是 \(x\) 的连续函数，且其绝对值不超过 \(L\)。固定 \(x\) 后，差商关于 \(t\ne0\) 也连续，因此 \(t\to0\) 的上下极限可只在非零有理数上计算。于是，方向导数存在集是两个 Borel 上下极限相等的集合，因而为 Borel 集。在此集合上记 \(g_v=D_vf\)，其外取 \(g_v=0\)，便得到满足 \(|g_v|\leq L\) 的 Borel 函数。

沿每条平行于 \(v\) 的直线，函数 \(t\mapsto f(y+tv)\) 都是 \(L\)-Lipschitz 的。在逐个有界区间上应用\cref{b1:thm:one-dimensional-lipschitz-ftc}，再取可数并，说明它在这条直线上几乎处处可导。若 \(n\geq2\)，选择以 \(v\) 为第一轴的正交坐标，由\cref{b1:thm:invertible-linear-change}及 Fubini 定理可知，整个 \(\mathbb R^n\) 中方向导数不存在的 Borel 集也为零测集。这里可先对每个有界坐标箱应用 Tonelli 定理，再穷竭全空间。当 \(n=1\) 时直接使用一维结论即可。至此，对每个固定的 \(v\)，都有
\[
 \Delta_{v,t}f\longrightarrow g_v
 \quad\text{几乎处处},\quad t\to0.
\]

现在取任意 \(\phi\in C_c^1(\mathbb R^n)\)。平移换元给出恒等式
\[
 \int_{\mathbb R^n}\Delta_{v,t}f(x)\cdot\phi(x)\,\mathrm dx
 =
 \int_{\mathbb R^n}f(x)
             \frac{\phi(x-tv)-\phi(x)}t\,\mathrm dx.
\]
左端的被积函数受 \(L|\phi|\) 支配。右端在 \(|t|\leq1\) 时仅涉及一个固定的紧集，\(f\) 在其上有界；而 \(\phi\) 的差商由 \(\|\nabla\phi\|_\infty\) 控制。因此两边都可以应用控制收敛定理，得到
\begin{equation}\label{b1:eq:directional-translation-test}
 \int g_v\phi=-\int fD_v\phi.
\end{equation}
对坐标方向 \(e_i\) 记 \(g_i=g_{e_i}\)。测试函数的普通微分给出 \(D_v\phi=\sum_{i=1}^n v_i\partial_i\phi\)，故再次使用\eqref{b1:eq:directional-translation-test}便有
\[
 \int\left(g_v-\sum_{i=1}^n v_ig_i\right)\phi=0.
\]
由\cref{b1:lem:compact-tests-separate}，
\[
 g_v=\sum_{i=1}^n v_ig_i
 \quad\text{几乎处处}.
\]
这条等式目前只对每个固定方向分别成立；不能直接合并不可数多个方向的例外集。

在单位球面 \(S^{n-1}\) 中取一个可数稠密集 \(V\)，并使它包含全部坐标方向。把这些方向中导数不存在或上述线性关系不成立的零测集取并，得到一个零测集 \(N\)。对每个 \(x\notin N\)，令 \(g(x)=(g_1(x),\ldots,g_n(x))\)，则同时有
\[
 D_vf(x)=g(x)\cdot v\quad(v\in V).
\]
又因 \(|D_vf(x)|\leq L\)，由 \(V\) 的稠密性和线性函数 \(v\mapsto g(x)\cdot v\) 的连续性可知 \(|g(x)|\leq L\)。

固定这样的 \(x\)，并对 \(v\in S^{n-1}\)、\(t>0\) 定义
\[
 Q(v,t)=\frac{f(x+tv)-f(x)}t-g(x)\cdot v.
\]
Lipschitz 条件给出一个与 \(t\) 无关的方向误差界：
\[
 |Q(v,t)-Q(w,t)|
 \leq\bigl(L+|g(x)|\bigr)|v-w|
 \leq2L|v-w|.
\]
给定 \(\varepsilon>0\)，从 \(V\) 中取单位球面的一个有限 \(\varepsilon/\bigl(2(1+2L)\bigr)\)-网，记其方向为 \(v_1,\ldots,v_k\)。在这有限个方向上分别使用 \(Q(v_j,t)\to0\)，可以选择共同的 \(\delta>0\)，使
\[
 |Q(v_j,t)|<\varepsilon/2
 \quad(1\leq j\leq k,\ 0<t<\delta).
\]
对任意单位方向 \(v\)，选取网中与它足够接近的 \(v_j\)，再用方向误差界，得到 \(|Q(v,t)|<\varepsilon\)。任意非零增量都可写成 \(h=|h|v\)，因此
\[
 \lim_{h\to0}\frac{|f(x+h)-f(x)-g(x)\cdot h|}{|h|}=0.
\]
这就证明了全空间实值情形。

若值域为 \(\mathbb R^m\)，对有限个分量分别应用实值结论，并删去这些零测集的并。各分量的余项除以 \(|h|\) 都趋零，有限维欧氏范数也就趋零，所得微分矩阵的各行由相应的分量微分给出。

最后处理局部情形。由局部 Lipschitz 条件和欧氏空间的可数拓扑基，可用可数个开球覆盖 \(U\)，使 \(f\) 在每个球上具有一个有限 Lipschitz 常数。对每个球和每个实分量，用\cref{b1:thm:mcshane-extension}延拓到全空间，再应用已经证明的结论。延拓函数在球内与原分量局部相同，所以在那里具有相同的可微性。可数个球及有限个分量的例外零测集之并仍为零测集。算子范数界则直接来自原函数的局部 Lipschitz 不等式，如上一小节所示。
\end{proof}

证明中的有限方向网是全微分出现的关键。平移测试先在一个共同的大集合上把方向导数确定为线性函数，Lipschitz 估计再使不同方向之间的差商不会失去控制；两者缺一，都不能从逐方向极限推出任意增量下的线性逼近。

\subsection{微分的 Borel 代表与大集上的一致线性化}
\label{b1:subsec:160204}

后续对微分矩阵作积分或划分水平集时，需要一个明确的可测代表。只说“微分几乎处处存在”还没有规定异常点上的取值。下面把真正的可微点集合与任意补值区别开来。

\begin{proposition}\label{b1:prop:lipschitz-differentiability-borel}
设 \(U\subseteq\mathbb R^n\) 为开集，\(f:U\rightarrow\mathbb R^m\) 局部 Lipschitz。它的可微点集合 \(\mathcal D_f\) 为 Borel 集。约定矩阵函数 \(Df\) 在 \(\mathcal D_f\) 上取真实全微分，在 \(U\setminus\mathcal D_f\) 上取零矩阵，则 \(Df:U\rightarrow\mathbb R^{m\times n}\) 为 Borel 可测函数。
\end{proposition}
\begin{proof}
对每个坐标方向 \(e_i\)，考虑 Borel 函数列
\[
 q_{i,k}(x)=
 \begin{cases}
 2^k\bigl(f(x+2^{-k}e_i)-f(x)\bigr),&x+2^{-k}e_i\in U,\\
 0,&x+2^{-k}e_i\notin U.
 \end{cases}
\]
对固定 \(x\in U\)，第一种情形最终总成立。由可数 Cauchy 条件，所有 \(n\) 列都具有有限极限的点组成 Borel 集 \(C\)。在 \(C\) 上以这些极限为列组成矩阵 \(A(x)\)，在 \(U\setminus C\) 上置零，则 \(A\) 为 Borel 矩阵函数。每个真正的可微点都在 \(C\) 中，且此时 \(A(x)=Df(x)\)。

为了检验全部增量，而不只是这几列差商，令
\[
 r_j(x)=
 \sup_{\substack{q\in\mathbb Q^n,\ 0<|q|<1/j\\x+q\in U}}
       \frac{|f(x+q)-f(x)-A(x)q|}{|q|}.
\]
对固定的非零有理向量 \(q\)，商在 \(x+q\in U\) 的 Borel 区域上可测，在其外补为零；可数上确界因而给出非负扩展实值 Borel 函数 \(r_j\)。固定 \(x\) 后，商关于允许的非零增量连续，所以有理增量的上确界等于全部实增量的上确界。于是
\[
 \mathcal D_f=C\cap\{\,x:\lim_{j\to\infty}r_j(x)=0\,\}.
\]
右端为 Borel 集，并与可微性定义完全相符：从左向右用真实微分作为候选矩阵，从右向左则由全部增量的误差趋零得到微分 \(A(x)\)。最后把 \(A\) 限于 \(\mathcal D_f\)，其外取零，就得到所述 Borel 代表。
\end{proof}

本书在后文使用 \(Df\) 的可测性时沿用这个约定。异常点上的零矩阵是补值，不表示函数在那里真的具有零微分。由 Rademacher 定理，这些补值只发生在零测集上；若 \(f\) 为全局 \(L\)-Lipschitz，则这个代表处处满足 \(\|Df\|\leq L\)。

\ProseParagraphBreak

逐点的余项估计还可以在删去小测度集合后统一起来。

\begin{proposition}\label{b1:prop:uniform-linearization-large-set}
设 \(f:\mathbb R^n\rightarrow\mathbb R^m\) 为 \(L\)-Lipschitz 映射，\(E\subseteq\mathbb R^n\) 为有限测度 Lebesgue 可测集。对每个 \(\varepsilon>0\)，存在紧集 \(K\subseteq E\cap\mathcal D_f\)，满足 \(m_n(E\setminus K)<\varepsilon\)，并存在非负函数 \(\omega(r)\to0\)（\(r\downarrow0\)），使
\[
 |f(x+h)-f(x)-Df(x)h|
 \leq\omega(|h|)|h|
 \quad(x\in K,\ h\ne0).
\]
\end{proposition}
\begin{proof}
在上一个证明的 \(r_j\) 中改用已经约定的 Borel 矩阵 \(Df\)。所得函数仍可测，并处处不超过 \(2L\)；在 \(\mathcal D_f\) 上有 \(r_j\downarrow0\)。对有限测度集 \(E\cap\mathcal D_f\) 应用\cref{b1:cor:egoroff-compact}，得到所需紧集 \(K\)，使 \(r_j\to0\) 在 \(K\) 上一致成立。由于 \(E\setminus\mathcal D_f\) 为零测集，可以同时保证 \(m_n(E\setminus K)<\varepsilon\)。

若 \(K\) 为空，取 \(\omega=0\)。否则定义
\[
 \omega(r)=
 \sup_{\substack{x\in K\\0<|h|\leq r}}
       \frac{|f(x+h)-f(x)-Df(x)h|}{|h|}.
\]
它不超过 \(2L\)，并且在 \(r<1/j\) 时不超过 \(\sup_{x\in K}r_j(x)\)。因此 \(\omega(r)\to0\)，所需误差界由定义直接得到。
\end{proof}

这个结论只统一了保留点上的一阶余项，并没有构造一个在整个空间上连续可微的延拓。下一节研究近似微分时，将允许增量估计忽略密度为零的方向区域；本节的真实微分及其 Borel 代表，则提供与这种较弱局部性质比较的基准。

% !TeX root = ../../Book1.tex
\section{近似微分}
\label{b1:sec:1603}

Rademacher 定理对 Lipschitz 映射给出通常意义的微分。
一般可测函数却会因零测集上的改值而失去连续性，这时普通微分并不适合作为几乎处处性质的唯一表述。
近似微分允许忽略在观察点附近所占比例趋于零的部分。
本节先比较两种微分，再证明一个分片结论：几乎处处近似可微的有限值可测映射，
除去零测集以后，可以限制在可数个紧集上分别成为 Lipschitz 映射。

% 来源范围：近似微分定义及唯一性沿用本书140504；
% 由密度一集合恢复普通极限用1405已证性质。
% Fremlin2016Measure 262N实际读到的是通常相对可微映射的Lipschitz分片，
% 不能把它直接引作这里近似可微版本的证明。
% 本节用统一球内坏点比例、两球相交与共同好点完整证明所需分片结论。

\subsection{近似可微性}
\label{b1:subsec:160301}

沿用\cref{b1:def:approximate-differential}，以下 \(\Omega\subseteq\mathbb R^d\) 为开集，
\(f:\Omega\rightarrow\mathbb R^m\) 有限值且是 Lebesgue 可测的。
若 \(f\) 在 \(x\) 处近似可微，把其唯一的线性映射记为
\(\operatorname{ap}Df(x)\)。记 \(A=\operatorname{ap}Df(x)\)，则对每个 \(\varepsilon>0\)，
\begin{equation}\label{b1:eq:approximate-differential-bad-set}
 \lim_{r\downarrow0}\frac{
 m_d\bigl(\{\,y\in B(x,r)\setminus\{x\}:
 |f(y)-f(x)-A(y-x)|>\varepsilon|y-x|\,\}\bigr)}
 {m_d(B(x,r))}=0.
\end{equation}
半径充分小时球包含于 \(\Omega\)。这个定义把函数值 \(f(x)\) 本身固定在观察点上，
并没有先用近似极限另选中心值。
\symidx{\(\operatorname{ap}Df(x)\)：近似微分}

\ProseParagraphBreak

\subsecref{b1:subsec:140504}已经用锥形方向上的正体积比例证明唯一性：
一个非零线性映射的相对大小 \(|Ah|/|h|\)，不可能在几乎所有小方向上趋于零。
因此近似微分仍是一张确定的线性映射，而不是随所忽略的集合变化的一族矩阵。

\begin{proposition}\label{b1:prop:approximate-differential-density-set}
设 \(\Omega\subseteq\mathbb R^d\) 为开集，\(f:\Omega\to\mathbb R^m\) 是有限值 Lebesgue 可测映射，\(x\in\Omega\)。给定线性映射 \(A:\mathbb R^d\to\mathbb R^m\)。\(f\) 在 \(x\) 处近似可微且近似微分为 \(A\)，当且仅当存在可测集
\(E\subseteq\Omega\)，使 \(E\) 在 \(x\) 处的密度为 \(1\)，并且
\[
 \lim_{\substack{y\to x\\y\in E,\ y\ne x}}
 \frac{|f(y)-f(x)-A(y-x)|}{|y-x|}=0.
\]
\end{proposition}
\begin{proof}
将\cref{b1:prop:approximate-limit-density-one}用于可测的误差商，即得正向结论。
反向若在密度为 \(1\) 的 \(E\) 上有这个极限，
则对每个 \(\varepsilon>0\)，在充分小的球中，误差商大于 \(\varepsilon\) 的点
只能落在 \(E\) 的补集。因此坏点比例趋于零，得到定义。
\end{proof}

这里选出的 \(E\) 可以依赖 \(x\)。
不能据此立刻声称存在一个公共满测集，使 \(f\) 在其上每一点都有相同的相对可微性质；
把逐点选择整理成可数个公共部分，正是本节最后需要补充的工作。

\ProseParagraphBreak

若 \(f\) 与 \(g\) 几乎处处相等，并且在当前观察点 \(x\) 也满足 \(f(x)=g(x)\)，
则两者在 \(x\) 处的近似可微性及近似微分相同。
因为两种误差商只会在一个零测集上不同，这个集合在每个球中的质量都为零。
若只改了 \(x\) 处的值，却没有保持 \(f(x)=g(x)\)，则点态结论可能改变。
不过对于几乎处处结论，可以再除去改值的零测集，所以几乎处处近似可微确实是等价类的性质。

\subsection{与普通微分的关系}
\label{b1:subsec:160302}

通常可微蕴含近似可微，且两种微分相等。
其逆命题对一般连续函数也不成立：很窄的尖峰可以使差商沿某列点保持较大，
同时让这些坏点在小球中所占比例趋于零。
下面先完整给出这个模型，便于直接比较两种微分。
但如果函数已有局部 Lipschitz 控制，窄峰就不能任意陡峭，逆命题反而成立。

\begin{proposition}[稀疏的连续尖峰]\label{b1:eg:sparse-peaks-approximate}
对正整数 \(n\)，取 \(a_n=2^{-n}\)、\(b_n=a_n^2/8\)，定义实值函数 \(f:\mathbb R\to\mathbb R\)
\[
 f(x)=\sum_{n\ge1}a_n\max\left\{1-\frac{|x-a_n|}{b_n},0\right\},
 \qquad J_n=(a_n-b_n,a_n+b_n).
\]
中心取几何级数使峰只向原点聚集，峰高取 \(a_n\) 保留非零差商；
半宽取更小的二次量，则使峰的总支撑比例消失。由于 \(b_n\le a_n/16\)，
\[
 J_n\subset(15a_n/16,17a_n/16),
\]
这些区间两两分离。每个非零点附近只有有限个峰，故 \(f\) 在那里连续。
在 \(J_n\) 上有 \(0\le f(x)\le a_n\le16x/15\)，其余点函数为零；
因此 \(f(0)=0\) 且函数在原点也连续。
若 \(J_n\) 与 \((-r,r)\) 相交，则 \(a_n<2r\)，从而
\[
 m_1\left(\bigcup_nJ_n\cap(-r,r)\right)
 \le\frac14\sum_{a_n<2r}a_n^2\le\frac43r^2.
\]
除以 \(2r\) 后趋于零。对每个 \(\varepsilon>0\)，差商绝对值大于
\(\varepsilon\) 的点都在这些峰内，故 \(\operatorname{ap}Df(0)=0\)。
另一方面 \(f(a_n)/a_n=1\)，而 \(f(-a_n)/(-a_n)=0\)，普通导数不存在。
峰的斜率绝对值为 \(a_n/b_n=8/a_n\to\infty\)，所以它也不在原点的任何邻域内 Lipschitz。
\end{proposition}

\cref{b1:fig:visual-f11}画的是这个模型；坏点的体积比例与峰点差商
分别对应近似微分和普通微分的不同检验。

\begin{proposition}\label{b1:prop:lipschitz-approximate-ordinary}
设 \(\Omega\subseteq\mathbb R^d\) 为开集，\(f:\Omega\to\mathbb R^m\) 是有限值 Lebesgue 可测映射，\(x\in\Omega\)。若 \(f\) 在 \(x\) 的某邻域内 Lipschitz，并在 \(x\) 处近似可微，
则 \(f\) 在 \(x\) 处通常可微，且
\[
 Df(x)=\operatorname{ap}Df(x).
\]
\end{proposition}
\begin{proof}
记 \(A=\operatorname{ap}Df(x)\)，令
\[
 R(y)=f(y)-f(x)-A(y-x).
\]
在所取邻域内，\(R\) 的 Lipschitz 常数不超过 \(M=L+\|A\|\)，且 \(R(x)=0\)。
若 \(M=0\)，结论直接成立。设 \(M>0\)。
若 \(R(y)=o(|y-x|)\) 不成立，就存在 \(\varepsilon>0\) 及
\(y_j\to x\)，使 \(r_j=|y_j-x|>0\) 且 \(|R(y_j)|\geq\varepsilon r_j\)。
取 \(0<c<\min\{1/2,\varepsilon/(2M)\}\)。
对 \(z\in B(y_j,cr_j)\)，当 \(j\) 充分大时，Lipschitz 估计给出
\[
 |R(z)|\geq |R(y_j)|-M|z-y_j|>\frac{\varepsilon r_j}{2},
 \quad |z-x|<(1+c)r_j.
\]
所以这些点全满足
\[
 \frac{|R(z)|}{|z-x|}>\frac{\varepsilon}{2(1+c)}.
\]
而球 \(B(y_j,cr_j)\) 包含于 \(B(x,(1+c)r_j)\)，
在后一个球中所占体积比例恒为 \(\bigl(c/(1+c)\bigr)^d>0\)。
这与误差商近似趋于零矛盾。
\end{proof}

这个论证说明 Lipschitz 条件在这里的具体作用：
若某点的误差达到距离的一定比例，附近一整小球的误差也不能骤然消失。
所以一个例外方向无法独自存在，它会带出一块有固定相对体积的坏点区域。
近似极限恰好排除了这种情况。

\ProseParagraphBreak

近似微分也有直接的代数运算规则。
若 \(f,g\) 在同一点近似可微，则任意固定线性组合也近似可微，
其微分是微分的相应线性组合：用三角不等式后，新的坏集包含在两个旧坏集之并中即可。
设 \(V\subseteq\mathbb R^m\) 为开集，\(f:\Omega\to V\) 为有限值 Lebesgue 可测映射，
\(\Phi:V\to\mathbb R^\ell\) 为 Borel 可测映射。若 \(f\) 在 \(x\) 处近似可微，
\(\Phi\) 在 \(f(x)\) 处通常可微，则 \(\Phi\circ f\) 在 \(x\) 处近似可微，并且
\begin{equation}\label{b1:eq:approximate-chain-rule}
 \operatorname{ap}D(\Phi\circ f)(x)
 =D\Phi(f(x))\circ\operatorname{ap}Df(x).
\end{equation}
复合 \(\Phi\circ f\) 的 Lebesgue 可测性由 \(\Phi\) 的 Borel 可测性保证。
为验证此式，取前一命题之前得到的密度一集合 \(E\)，在其上有
\[
 f(y)-f(x)=A(y-x)+o(|y-x|).
\]
因而沿 \(E\) 有 \(|f(y)-f(x)|=O(|y-x|)\) 且 \(f(y)\to f(x)\)。
将它代入 \(\Phi\) 的普通可微展开，所得余项仍为 \(o(|y-x|)\)，
再由密度一集合判据得到链式法则。
此处假设外层 \(\Phi\) 通常可微，不能不加条件地换成仅近似可微：
内层映射可能把大块集合压到外层的异常集合上。

\ProseParagraphBreak

另一个常用结论是，在可测水平集的几乎处处，近似微分为零。
确切地说，若 \(E=\{x:f(x)=v\}\)，则在 \(E\) 的每个密度点，
沿 \(E\) 函数恒为 \(v=f(x)\)，故零映射满足密度一集合判据。
Lebesgue 密度定理说明 \(E\) 中几乎每点都是这样的点。
特别地，若 \(f\) 本身 Lipschitz，则在该水平集的几乎处处有 \(Df=0\)。

\subsection{近似可微映射的 Lipschitz 分片}
\label{b1:subsec:160303}

近似可微比 Lipschitz 条件弱得多，却仍能在可数个部分上恢复距离控制。
下面的证明不使用光滑延拓定理，而是从近似微分所给的小球坏点比例出发，
为两个中心找到一个共同的好点。

\begin{theorem}\label{b1:thm:approximate-lipschitz-pieces}
设 \(\Omega\subseteq\mathbb R^d\) 为开集，\(f:\Omega\rightarrow\mathbb R^m\)
为有限值 Lebesgue 可测映射。下列两条等价：
\begin{equivalences}
\item\label{b1:item:approximate-pieces-ae}映射 \(f\) 在 \(\Omega\) 中几乎处处近似可微。
\item\label{b1:item:approximate-pieces-compact}存在可数个紧集 \(K_j\subseteq\Omega\)，使
\[
 m_d\left(\Omega\setminus\bigcup_jK_j\right)=0,
\]
且每个限制 \(f|_{K_j}\) 都是 Lipschitz 映射，其常数可以依赖于 \(j\)。
\end{equivalences}
\end{theorem}
\begin{proof}
先证\itemref{b1:item:approximate-pieces-ae}蕴含%
\itemref{b1:item:approximate-pieces-compact}。
由于 \(\Omega\) 是开集，把 \(f\) 在 \(\mathbb R^d\setminus\Omega\) 上补零后，
所得映射仍为有限值 Lebesgue 可测映射；再由
\cref{b1:thm:lebesgue-borel-representative}选取一个有限值 Borel 代表。
原函数与代表只差一个 Lebesgue 零测集；最终选紧集时再避开一个包含该零测集的 Borel 零测集，
所得限制也就适用于原函数。以下暂把这个 Borel 代表仍记为 \(f\)。

固定 \(0<\eta<2^{-d-1}\)。对正整数 \(n\)，令 \(A_n\) 由所有满足以下条件的
\(x\in\Omega\) 组成：
\[
 \begin{gathered}
 |x|\leq n,\quad |f(x)|\leq n,\\
 m_d\bigl(\{\,y\in B(x,r):
       |f(y)-f(x)|>n|y-x|\,\}\bigr)
 \leq\eta\,\omega_d r^d
 \quad(0<r<1/n).
 \end{gathered}
\]
这些集合是 Borel 集。事实上，被积的坏点指标函数关于 \((x,y)\) 为 Borel 函数，
Tonelli 定理保证对 \(y\) 积分后关于 \(x\) 可测；
而半径条件只需对有理数 \(r\) 检查，任意半径由开球从内逼近及 \(r^d\) 连续得到。

若 \(f\) 在 \(x\) 处近似可微，其微分为 \(T\)，选整数 \(n_0>\|T\|+1\)。
当 \(|f(y)-f(x)|>n_0|y-x|\) 时，误差商必大于 \(1\)。
由近似微分，充分小球中的这类点所占比例小于 \(\eta\)。
再把 \(n_0\) 增大到某个 \(n\)，使 \(1/n\) 小于所需半径，
并使 \(|x|,|f(x)|\leq n\)，便有 \(x\in A_n\)。
增大阈值只会缩小坏集，所以这一步保持估计。
因此 \(\bigcup_n A_n\) 包含几乎所有点。

证明 \(f|_{A_n}\) 为 Lipschitz。取不同的 \(x,y\in A_n\)，记 \(\rho=|x-y|\)。
若 \(\rho<1/(2n)\)，考虑
\[
 B\left(\frac{x+y}{2},\rho\right)
 \subset B(x,2\rho)\cap B(y,2\rho).
\]
两个中心各自的坏点集，在其半径为 \(2\rho\) 的球内至多具有质量
\(\eta\omega_d(2\rho)^d\)。
二者质量之和严格小于 \(\omega_d\rho^d\)，所以在上述公共球中可找到一点 \(z\)，
它对两个中心都不是坏点。于是
\[
 |f(x)-f(y)|
 \leq|f(x)-f(z)|+|f(z)-f(y)|
 \leq n\bigl(|x-z|+|z-y|\bigr)
 \leq4n\rho.
\]
若 \(\rho\geq1/(2n)\)，则由 \(|f(x)|,|f(y)|\leq n\)，有
\[
 |f(x)-f(y)|\leq2n\leq4n^2\rho.
\]
合起来得到 \(f|_{A_n}\) 的 Lipschitz 常数不超过 \(4n^2\)。

每个 \(A_n\) 有界，Lebesgue 测度有限。
由紧内正则性，可以在 \(A_n\) 内选可数个紧集，其并除去 \(A_n\) 的一个零测集；
若原函数换过代表，就先从 \(A_n\) 删去前述 Borel 零测集。
把所有这些紧集列为一列 \((K_j)\)，得到所需覆盖及 Lipschitz 限制。

反过来，假设\itemref{b1:item:approximate-pieces-compact}成立。
由逐坐标 McShane 延拓，将 \(f|_{K_j}\) 延拓为全空间 Lipschitz 映射 \(F_j\)。
Rademacher 定理保证 \(F_j\) 几乎处处可微，Lebesgue 密度定理保证
\(K_j\) 中几乎每点都是 \(K_j\) 的密度点。
在同时满足两条件的 \(x\in K_j\)，沿 \(K_j\) 有 \(f=F_j\)，故
\[
 \lim_{\substack{y\to x\\y\in K_j,\ y\ne x}}
 \frac{|f(y)-f(x)-DF_j(x)(y-x)|}{|y-x|}=0.
\]
由密度一集合判据，\(f\) 在 \(x\) 处近似可微。
合并可数个零测例外并使用紧集族几乎处处覆盖，即得%
\itemref{b1:item:approximate-pieces-ae}。
\end{proof}

定理只保证分片上的 Lipschitz 性质，常数可能趋于无穷，片与片之间也不提供统一距离估计。
因此不能由它倒推出 \(f\) 局部 Lipschitz。
同样，紧片的像是紧的，可数个紧片像之并也是 Borel 集，
但被遗漏的定义域零测集可能在一般可测映射下有很大的像；
没有额外的零测集映射性质时，不能忽略这部分像。

\begin{corollary}\label{b1:cor:approximate-derivative-measurable}
设 \(\Omega\subseteq\mathbb R^d\) 为开集，\(f:\Omega\to\mathbb R^m\) 是有限值 Lebesgue 可测映射。若 \(f\) 几乎处处近似可微，则其近似微分的矩阵系数有 Borel 可测代表。
\end{corollary}
\begin{proof}
沿前一定理取紧集 \(K_j\) 和全空间延拓 \(F_j\)。
由\cref{b1:prop:lipschitz-differentiability-borel}，每个 \(DF_j\)
在可微集外补零后是 Borel 矩阵函数。
置 \(E_j=K_j\setminus\bigcup_{i<j}K_i\)，在 \(E_j\) 上取矩阵 \(DF_j\)，
在所有 \(E_j\) 之外取零，得到 Borel 矩阵函数。
前一定理的反向证明表明，这个函数在几乎处处的点等于 \(\operatorname{ap}Df\)。
\end{proof}

普通微分从相邻点的逐点行为产生，近似微分则先允许丢掉相对稀薄的异常点。
本节的分片定理说明，两者之间可以通过可数个 Lipschitz 限制建立联系。
后续面积公式将需要这种联系，但还必须另行处理像的测度与重数；
本节没有把这些未建立的结论包含进“分片”二字之中。

\cref{b1:eg:sparse-peaks-approximate}的支撑比例与峰点差商如\cref{b1:fig:visual-f11}所示；两者检验的量不能互换。
% F11：本书原创矢量示意；依据所在节的定义、证明或明确模型。
\begin{figure}[htbp]
\centering
\begin{tikzpicture}[x=1cm,y=1cm,>=stealth,every node/.style={font=\small},line width=.65pt]
\draw[->,gray] (-0.12,0)--(5.4,0) node[right] {\(x\)};\draw[->,gray] (0,-0.12)--(0,2.8);
\node[] at (2.4,3.4) {窄峰函数};
\draw[AnalysisBlue,thick] plot coordinates {(4.2188,0.0000) (4.5000,2.5000) (4.7812,0.0000)};
\draw[AnalysisBlue,thick] plot coordinates {(2.1797,0.0000) (2.2500,1.2500) (2.3203,0.0000)};
\draw[AnalysisBlue,thick] plot coordinates {(1.1074,0.0000) (1.1250,0.6250) (1.1426,0.0000)};
\draw[AnalysisBlue,thick] plot coordinates {(0.5581,0.0000) (0.5625,0.3125) (0.5669,0.0000)};
\draw[AnalysisBlue,thick] plot coordinates {(0.2802,0.0000) (0.2812,0.1562) (0.2823,0.0000)};
\draw[AnalysisBlue,thick] plot coordinates {(0.1404,0.0000) (0.1406,0.0781) (0.1409,0.0000)};
\draw[->,gray] (6.88,0)--(12.1,0) node[right] {\(x\)};\draw[->,gray] (7,-0.12)--(7,2.8);
\node[] at (9.3,3.4) {峰点差商不趋于零};
\fill[orange!80!black](11.5,2)circle(1.7pt);\fill[orange!80!black](9.25,2)circle(1.7pt);\fill[orange!80!black](8.125,2)circle(1.7pt);\fill[orange!80!black](7.5625,2)circle(1.7pt);\fill[orange!80!black](7.28125,2)circle(1.7pt);\fill[orange!80!black](7.140625,2)circle(1.7pt);\draw[gray,dashed] plot coordinates {(7.0000,2.0000) (11.8000,2.0000)};
\node[] at (9.3,-0.6) {\(f(a_n)/a_n=1\)，坏集密度趋零};

\end{tikzpicture}
\caption[稀疏尖峰与近似微分]{稀疏尖峰与近似微分。使用正文\cref{b1:eg:sparse-peaks-approximate}的参数 \(a_n=2^{-n}\)、\(b_n=a_n^2/8\)。峰中心和峰高均为 \(a_n\)，半宽为 \(b_n\)。左图画出前六个三角峰（坐标作示意缩放）；全部峰的支撑在原点密度为零，但峰点差商恒为一。}
\label{b1:fig:visual-f11}
\end{figure}


% !TeX root = ../../Book1.tex
\section{绝对连续函数与微积分基本定理}
\label{b1:sec:1604}

绝对连续性把“小区间族的总长度很小”转化为“函数增量的总和很小”。它比一致连续性强，又允许导数在局部无界；更重要的是，它恰好给出可以由可积导数恢复的实函数。本节的结果是后续一维 Sobolev 代表理论的直接前置，而不属于可以跳过的度量曲线选读内容。

% 实读来源：Fremlin2016Measure，作者官方 chap22.pdf，225B--E（该文件印刷 pp.27--29）。
% 积分表示以变差、单调分解、广义逆推前和 RN 补写，不把未写的 Stieltjes 构造当作先修。

\subsection{绝对连续性与全变差}
\label{b1:subsec:160401}

\begin{definition}
\label{b1:def:absolutely-continuous-function}
设 \(f:[a,b]\to\mathbb R\)，其中 \(-\infty<a<b<\infty\)。若对每个 \(\varepsilon>0\)，存在 \(\delta>0\)，使任意有限族内点互不相交的区间 \([s_i,t_i]\subseteq[a,b]\) 满足
\[
\sum_i(t_i-s_i)<\delta
\quad\Longrightarrow\quad
\sum_i\bigl|f(t_i)-f(s_i)\bigr|<\varepsilon,
\]
则称 \(f\) 为\term[juedui-lianxu-hanshu]{绝对连续函数}[absolutely continuous function]。
\end{definition}

这里只取一个区间，便可看出绝对连续函数一致连续。Lipschitz 函数满足定义，因为增量绝对值之和不超过 Lipschitz 常数乘以区间长度之和；反向一般不成立，例如本节末的平方根函数。

\begin{lemma}
\label{b1:lem:ac-real-variation}
设 \(-\infty<a<b<\infty\)。设 \(f:[a,b]\to\mathbb R\) 绝对连续，定义
\[
V(t)=\sup\Bigl\{\sum_{i=1}^n
  \bigl|f(t_i)-f(t_{i-1})\bigr|:
a=t_0<\cdots<t_n=t\Bigr\},
\quad V(a)=0.
\]
则 \(V\) 有限、单调不减且绝对连续。对 \(a\leq s<t\leq b\)，有
\[
\bigl|f(t)-f(s)\bigr|\leq V(t)-V(s),
\]
而右端正是 \(f\) 在子区间 \([s,t]\) 上的全变差。
\end{lemma}
\begin{proof}
在一个分割中插入中间点，由三角不等式，增量绝对值之和不会减小；反过来，可以拼接两个子区间的分割。这证明分割上确界对子区间可加，允许暂取无穷值。

对 \(\varepsilon=1\) 取绝对连续性定义中的 \(\delta\)，把 \([a,b]\) 分成有限个长度小于 \(\delta\) 的区间。在其中任一区间内，任意分割的各小段总长度小于 \(\delta\)，所以增量绝对值之和小于 \(1\)。把一个任意分割加上这些预定端点，便得到总上界，故 \(V(b)<\infty\)。单调性、子区间公式及增量估计随之成立。

最后，给定 \(\varepsilon>0\)，对绝对连续性取使增量绝对值之和小于 \(\varepsilon/2\) 的 \(\delta\)。若一族不交区间的总长度小于 \(\delta\)，在每个区间内部选分割来逼近其分割上确界。所有分割小段仍组成总长度小于 \(\delta\) 的不交族。因此，取这些有限个上确界后，
\[
\sum_i\bigl(V(t_i)-V(s_i)\bigr)\leq\varepsilon/2<\varepsilon.
\]
这正是 \(V\) 的绝对连续性。
\end{proof}

\subsection{微积分基本定理}
\label{b1:subsec:160402}

Fremlin 的《Measure Theory》\cite[225B--E]{Fremlin2016Measure}系统讨论绝对连续函数及其积分表示。下面采用\chapref{b1:ch:06}的 Radon--Nikodym 定理来证明表示方向，使这里不依赖尚未建立的一般有界变差函数微分理论。

\begin{theorem}[微积分基本定理，绝对连续函数]
\label{b1:thm:ac-fundamental-calculus}
设 \(-\infty<a<b<\infty\)。实函数 \(f:[a,b]\to\mathbb R\) 绝对连续，当且仅当存在 \(g\in L^1([a,b])\)，使
\[
f(t)=f(a)+\int_a^t g(r)\,\mathrm dr
\quad(t\in[a,b]).
\]
在这种情形下，\(f'=g\) 几乎处处，因而可以取 \(g=f'\)。
\end{theorem}
证明的正向先以全变差函数把一般 \(f\) 分解为两个单调绝对连续函数：
\[
 f\longrightarrow V\longrightarrow P,Q.
\]
对其中任一单调函数 \(F\)，再用广义逆 \(q\) 把区间上的 Lebesgue 测度推前为增量测度
\(\mu_F\)，由绝对连续性证明 \(\mu_F\ll m_1\)，最后应用 Radon--Nikodym 定理取得积分密度。
反向则直接用积分的绝对连续性。因而正向证明的主链是
\[
 F\longrightarrow q\longrightarrow\mu_F
 \longrightarrow\mu_F\ll m_1
 \longrightarrow\frac{\mathrm d\mu_F}{\mathrm dm_1}.
\]
\begin{proof}
先设 \(f\) 绝对连续，令 \(V\) 为\cref{b1:lem:ac-real-variation}中的全变差函数。两个函数
\[
P(t)=\frac{V(t)+f(t)-f(a)}2,
\quad
Q(t)=\frac{V(t)-f(t)+f(a)}2
\]
都单调不减，因为 \(V(t)-V(s)\geq|f(t)-f(s)|\)。它们还绝对连续，在 \(a\) 处为零，且 \(f-f(a)=P-Q\)。所以只需为任意单调不减的绝对连续函数 \(F\)，在 \(F(a)=0\) 的归一化下，构造积分密度。

令 \(M=F(b)\)。若 \(M=0\)，则 \(F=0\)。若 \(M>0\)，对 \(0<u<M\) 定义
\[
q(u)=\min\{t\in[a,b]:F(t)\geq u\}.
\]
这个最小值由 \(F\) 的连续性与单调性存在，并且
\[
q(u)\leq t\quad\Longleftrightarrow\quad u\leq F(t).
\]
所以 \(q\) 是 Borel 可测函数。把 \((0,M)\) 上的 Lebesgue 测度经 \(q\) 推前，得到 \([a,b]\) 上的有限正 Borel 测度 \(\mu_F\)，满足
\[
\mu_F([a,t])=F(t),
\quad
\mu_F((s,t])=F(t)-F(s).
\]
又由 \(F\) 连续，这个测度没有原子，包括两个端点。

现证明 \(\mu_F\ll m_1\)。设 Borel 集 \(N\subseteq[a,b]\) 的 Lebesgue 测度为零，给定 \(\varepsilon>0\)，取 \(F\) 的绝对连续性所对应的 \(\delta\)。由 Lebesgue 测度的外正则性，存在开集 \(O\supseteq N\)，使 \(m_1(O)<\delta\)。开集 \(O\) 与 \([a,b]\) 的交可分成可数个内点互不相交的区间；端点有没有包含不影响 \(\mu_F\) 的值。每个有限子族的 \(\mu_F\) 测度之和，就是相应的 \(F\) 增量之和，不超过 \(\varepsilon\)。由可数可加性，\(\mu_F(N)\leq\mu_F(O\cap[a,b])\leq\varepsilon\)。令 \(\varepsilon\downarrow0\)，得到绝对连续性。

由\cref{b1:thm:radon-nikodym-finite}，存在非负 \(g_F\in L^1([a,b])\)，使 \(\mu_F=g_Fm_1\)，因此 \(F(t)=\int_a^t g_F\)。分别用于 \(P,Q\)，便得到所需 \(g=g_P-g_Q\)。这同时补出了本证明所需的单调函数生成测度的步骤。

反过来，若 \(f(t)=f(a)+\int_a^t g\)，则对任意不交区间族，
\[
\sum_i|f(t_i)-f(s_i)|
\leq\int_{\bigcup_i[s_i,t_i]}|g(r)|\,\mathrm dr.
\]
积分的绝对连续性，即\cref{b1:prop:integral-absolute-continuity}，证明 \(f\) 绝对连续。最后，在 \(g\) 的每个内部 Lebesgue 点 \(t\)，
\[
\left|\frac{f(t+h)-f(t)}h-g(t)\right|
\leq\frac1{|h|}\int_{t-|h|}^{t+|h|}|g(r)-g(t)|\,\mathrm dr
\longrightarrow0.
\]
\cref{b1:thm:lebesgue-point}遂给出 \(f'=g\) 几乎处处。
\end{proof}

例如，\(t\mapsto\sqrt t\) 在 \([0,1]\) 上绝对连续，因为它是 \((2\sqrt t)^{-1}\) 的积分，但它在零点附近不是 Lipschitz 的。另一方向的边界由\cref{b1:eg:cantor-function}中的 Cantor 函数给出：它连续、单调且导数几乎处处为零，却不是常数，所以不是绝对连续函数。仅有连续性和几乎处处可微性，还不足以保证微积分基本定理。

有限区间上的正则性条件决定导数能否恢复全部函数增量，见\cref{b1:tab:visual-t09}。
% Source fingerprint: 89a41d06e175d9aa5ca92aa3ddd8667a41bbec1369852f31d74b225a69b2e5b8
% T09: raw TeX cells preserved; no numeric highlighting.
\begin{table}[htbp]
\centering\small\renewcommand{\arraystretch}{1.18}\setlength{\tabcolsep}{4pt}
\caption[Lipschitz、绝对连续与有界变差]{Lipschitz、绝对连续与有界变差。这里的一维变差按划分定义。第27章从分布导数测度重新组织该理论；函数代表与端点取值始终按正文约定。}
\label{b1:tab:visual-t09}
\begin{tabularx}{\linewidth}{@{}>{\raggedright\arraybackslash}p{3.1cm}>{\raggedright\arraybackslash}p{4.7cm}>{\raggedright\arraybackslash}X@{}}
\toprule
有限区间上的函数类 & 导数与积分恢复 & 不能逆推的例子 \\
\midrule
Lipschitz & 绝对连续；导数几乎处处存在且本性有界；积分恢复函数增量 & 绝对连续函数不必 Lipschitz，如平方根 \\
绝对连续 & 导数属于 \(L^1\)，且 \(f(t)-f(a)=\int_a^t f^{\prime}(s)\,\mathrm{d}s\) & 仅有连续性和几乎处处可微不足以保证积分恢复 \\
有界变差 & 划分增量绝对值之和有统一上界；导数可积，但可有奇异变化 & Cantor 函数连续、几乎处处导数为零，却不是常值 \\
一般连续函数 & 不保证有界变差或几乎处处可微 & 本章三角峰例子可使划分增量和发散 \\
\bottomrule
\end{tabularx}
\end{table}


% !TeX root = ../../Book1.tex
\OptionalSection{绝对连续曲线、度量导数与长度}{b1:sec:1605}

在欧氏空间中，曲线的导数同时记录速度的大小与方向。一般度量空间没有向量差，方向未必能够定义，但相邻时刻的位置仍有距离。用距离除以时间差，可以保留速度的大小；要使它与长度的积分相联系，合适的条件是绝对连续性。

\ProseParagraphBreak

本节的参数区间固定为 \([a,b]\)，其中 \(a<b\)。空间只假设为度量空间，不要求线性、完备或局部紧。上一节已经建立实值绝对连续函数的积分表示；这里把一般曲线转换成可数个实值距离函数，使测度论微积分能够用于没有坐标的空间。

% 度量空间绝对连续性定义对照 Fremlin2016Measure 作者官方 chap26.tex 的 264Yp。
% 度量导数对照 Rossi--Mielke--Savare (2008)，A metric approach to a class of doubly
% nonlinear evolution equations and applications，作者公开版 §2.1，Proposition 2.2 及证明，pp.10--12：
% https://riccarda-rossi.unibs.it/pubblicazioni/Published%20papers/pubblicazioni-files/rossi-mielke-savare.pdf
% 正式发表信息：Ann. Scuola Norm. Sup. Pisa Cl. Sci. (5), 7(1), 97--169。
% 该文处理更广的非对称距离；本节只证对称度量，用可数距离测试和长度函数独立组织证明。

\subsection{绝对连续曲线}
\label{b1:subsec:160501}

\begin{definition}
\label{b1:def:absolutely-continuous-curve}
设 \((X,\mathrm{d})\) 为度量空间，\(\gamma:[a,b]\to X\)。若对每个 \(\varepsilon>0\)，存在 \(\delta>0\)，使任意有限族内点互不相交的区间 \([s_i,t_i]\subseteq[a,b]\) 满足
\[
\sum_i(t_i-s_i)<\delta
\quad\Longrightarrow\quad
\sum_i \mathrm{d}\bigl(\gamma(s_i),\gamma(t_i)\bigr)<\varepsilon,
\]
则称 \(\gamma\) 为\term[juedui-lianxu-quxian]{绝对连续曲线}[absolutely continuous curve]，记作 \(\gamma\in\operatorname{AC}([a,b];X)\)。当 \(X=\mathbb R\) 时，这一定义与上一节的绝对连续实值函数一致。
\end{definition}

这里只取一个区间，便可看出绝对连续曲线一致连续；它的像是紧的，因而是可分度量空间。Lipschitz 曲线满足定义，因为所有距离之和不超过 Lipschitz 常数乘以区间长度之和。绝对连续性允许不同位置具有不同的速度控制，但要求小的总时间长度只能产生小的总位移。

\begin{lemma}
\label{b1:lem:ac-variation}
设 \(\gamma\in\operatorname{AC}([a,b];X)\)，定义
\[
V(t)=\sup\Bigl\{\sum_{i=1}^n \mathrm{d}\bigl(\gamma(t_i),\gamma(t_{i-1})\bigr):
a=t_0<\cdots<t_n=t\Bigr\},
\quad V(a)=0.
\]
则 \(V\) 有限、单调不减且绝对连续。对 \(a\leq s<t\leq b\)，有
\[
\mathrm{d}\bigl(\gamma(s),\gamma(t)\bigr)\leq V(t)-V(s),
\]
而右端正是同一个分割上确界在子区间 \([s,t]\) 上的值。
\end{lemma}
\begin{proof}
在一个分割中插入中间点，由三角不等式，距离和不会减小；反过来，可以拼接两个子区间的分割。这证明分割上确界对子区间可加，允许暂取无穷值。

对 \(\varepsilon=1\) 取定义中的 \(\delta\)，把 \([a,b]\) 分成有限个长度小于 \(\delta\) 的区间。在其中任一区间内，任意分割的各小段总长度小于 \(\delta\)，所以距离和小于 \(1\)。把一个任意分割加上这些预定端点，便得到总上界，故 \(V(b)<\infty\)。单调性、子区间公式及位移估计随之成立。

最后，给定 \(\varepsilon>0\)，对绝对连续性取使距离和小于 \(\varepsilon/2\) 的 \(\delta\)。若一族不交区间的总长度小于 \(\delta\)，在每个区间内部选分割来逼近其分割上确界。所有分割小段仍组成总长度小于 \(\delta\) 的不交族。因此，取这些有限个上确界后，
\[
\sum_i\bigl(V(t_i)-V(s_i)\bigr)\leq\varepsilon/2<\varepsilon.
\]
这正是 \(V\) 的绝对连续性。
\end{proof}

由\cref{b1:lem:ac-variation}与\cref{b1:thm:ac-fundamental-calculus}立即得到一个等价而实用的条件：\(\gamma\) 绝对连续，当且仅当存在非负 \(g\in L^1([a,b])\)，使
\[
\mathrm{d}\bigl(\gamma(s),\gamma(t)\bigr)
\leq\int_s^t g(r)\,\mathrm dr
\quad(a\leq s<t\leq b).
\]
正向由\cref{b1:lem:ac-variation}知 \(V\) 绝对连续，故微积分基本定理给出
\(V(t)-V(s)=\int_s^tV'(r)\,\mathrm dr\)，并且 \(V'\geq0\) 几乎处处；取 \(g=V'\) 即可。
反向对不交区间求和，再用积分绝对连续性。于是绝对连续曲线确有一个可积的速度上界。下一小节将证明，这些上界中有一个几乎处处最小的选择。

\subsection{度量导数}
\label{b1:subsec:160502}

\begin{definition}
\label{b1:def:metric-derivative}
设 \(\gamma:[a,b]\to X\) 为曲线。若在 \(t\in(a,b)\) 处极限
\[
|\gamma'|_{\mathrm d}(t)
\coloneqq\lim_{h\to0}\frac{\mathrm{d}\bigl(\gamma(t+h),\gamma(t)\bigr)}{|h|}
\]
存在且有限，则称其为 \(\gamma\) 在 \(t\) 处的\term[duliang-daoshu]{度量导数}[metric derivative]，也称该处的度量速度。这里 \(|\gamma'|_{\mathrm d}\) 是一个整体记号，并不预设 \(\gamma'\) 作为向量存在。
\symidx{\(\lvert\gamma'\rvert_{\mathrm d}(t)\)：曲线的度量导数}
\end{definition}

度量导数及其最小控制性质，可以与 Rossi、Mielke、Savar\'e 的《A metric approach to a class of doubly nonlinear evolution equations and applications》\cite[Section 2.1, Proposition 2.2]{Rossi2008Metric}对照阅读。该文还允许非对称距离，服务于演化方程的度量表述；本节只讨论通常的对称度量，并将证明组织为距离函数与长度函数之间的比较。

\begin{theorem}[度量导数的存在性]
\label{b1:thm:metric-derivative}
若 \(\gamma\in\operatorname{AC}([a,b];X)\)，则其度量导数几乎处处存在。在例外零集上把它赋值为零后，它是可积函数，并且
\[
\mathrm{d}\bigl(\gamma(s),\gamma(t)\bigr)
\leq\int_s^t|\gamma'|_{\mathrm d}(r)\,\mathrm dr
\quad(a\leq s<t\leq b).
\]
若非负 \(g\in L^1([a,b])\) 也满足这个位移控制不等式，则 \(|\gamma'|_{\mathrm d}\leq g\) 几乎处处。
\end{theorem}
\begin{proof}
在紧像 \(\gamma([a,b])\) 中取可数稠密集 \(\{y_j:j\geq1\}\)，令
\[
f_j(t)=\mathrm{d}\bigl(\gamma(t),y_j\bigr).
\]
三角不等式给出 \(|f_j(t)-f_j(s)|\leq \mathrm{d}\bigl(\gamma(t),\gamma(s)\bigr)\)，因此每个 \(f_j\) 都是实值绝对连续函数。它们还完整记录像中两点的距离：
\[
\mathrm{d}(x,y)=\sup_j|\mathrm{d}(x,y_j)-\mathrm{d}(y,y_j)|
\quad(x,y\in\gamma([a,b])).
\]
其中“\(\geq\)”来自三角不等式；为得反向不等式，只需选 \(y_j\to x\)。

在全部 \(f_j\) 及引理中的 \(V\) 都可微的共同满测度集上，定义
\[
v(t)=\sup_j|f_j'(t)|,
\]
在其补集上赋值为零。导数可由可测差商的极限得到，因此 \(v\) 可测。由 \(|f_j(t)-f_j(s)|\leq|V(t)-V(s)|\)，有 \(0\leq v\leq V'\) 几乎处处，故 \(v\in L^1\)。实值微积分基本定理又给出
\[
|f_j(t)-f_j(s)|\leq\int_s^t v(r)\,\mathrm dr.
\]
对 \(j\) 取上确界，得到 \(\mathrm{d}\bigl(\gamma(s),\gamma(t)\bigr)\leq\int_s^tv\)。在 \([s,t]\) 的任意分割上求和并取上确界，进一步得到
\[
V(t)-V(s)\leq\int_s^t v(r)\,\mathrm dr.
\]
在 \(V\) 可微且 \(v\) 为 Lebesgue 点之处取差商，可知 \(V'\leq v\) 几乎处处。因此
\[
v=V'\quad\text{几乎处处}.
\]

现在取一个满足上述等式、且全部 \(f_j\) 和 \(V\) 可微的内部点 \(t\)。对每个 \(j\)，距离商不小于 \(|f_j(t+h)-f_j(t)|/|h|\)，所以
\[
\liminf_{h\to0}\frac{\mathrm{d}\bigl(\gamma(t+h),\gamma(t)\bigr)}{|h|}
\geq\sup_j|f_j'(t)|=v(t).
\]
另一方面，由\cref{b1:lem:ac-variation}，同一个距离商不超过 \(|V(t+h)-V(t)|/|h|\)，故其上极限不超过 \(V'(t)=v(t)\)。这证明极限存在，且 \(|\gamma'|_{\mathrm d}=V'=v\) 几乎处处，并给出可积性及位移估计。

若 \(g\) 是另一个可积控制函数，则在度量导数存在且 \(g\) 为 Lebesgue 点的地方，对 \(\gamma(t+h)\) 与 \(\gamma(t)\) 使用控制不等式，再除以 \(|h|\) 并令 \(h\to0\)，便得到 \(|\gamma'|_{\mathrm d}(t)\leq g(t)\)。
\end{proof}

这个证明的可数性来自曲线的紧像，而不是来自环境空间的可分性。它也没有把可数个实导数组合成一个向量导数，只是用这些导数确定了一个非负速度。

\ProseParagraphBreak

这种区别在无限维空间中确实会发生。取 \(X=L^1(0,1)\)，令 \(\gamma(t)=\mathbf1_{[0,t]}\)。则
\[
\|\gamma(t)-\gamma(s)\|_1=|t-s|,
\]
所以曲线是等距的，\(|\gamma'|_{\mathrm d}(t)=1\) 在所有内部点成立。然而，对充分小的 \(h>0\)，正向差商为 \(q_h=h^{-1}\cdot\mathbf1_{(t,t+h]}\)，并且
\[
\|q_h-q_{h/2}\|_1
=\int_t^{t+h/2}\frac1h\,\mathrm dr
+\int_{t+h/2}^{t+h}\frac1h\,\mathrm dr=1.
\]
差商不是 Cauchy 族，故在任何内部点都不存在 \(L^1\) 范数意义下的向量导数。一般 Banach 空间中的绝对连续曲线不能直接套用有限维 Rademacher 定理；度量导数的存在性却不受这个障碍影响。

\subsection{长度公式}
\label{b1:subsec:160503}

\begin{definition}
\label{b1:def:curve-length}
给定连续曲线 \(\gamma:[a,b]\to X\)，称
\[
\ell(\gamma)
\coloneqq\sup\Bigl\{\sum_{i=1}^n \mathrm{d}\bigl(\gamma(t_i),\gamma(t_{i-1})\bigr):
a=t_0<\cdots<t_n=b\Bigr\}
\]
为曲线的\term[quxian-changdu]{长度}[length of a curve]。若 \(\ell(\gamma)<\infty\)，则称其为\term[keqiuchang-quxian]{可求长曲线}[rectifiable curve]。
\symidx{\(\ell(\gamma)\)：参数化曲线的长度}
\end{definition}

长度计入曲线的行进过程，不只是像集的大小。绝对连续曲线的长度就是先前的 \(V(b)\)，而 \(V(t)\) 记录截至时刻 \(t\) 的累积长度。

\begin{theorem}[绝对连续曲线的长度公式]
\label{b1:thm:metric-curve-length}
设 \(\gamma\in\operatorname{AC}([a,b];X)\)。则对任意 \([s,t]\subseteq[a,b]\)，
\[
\ell\bigl(\gamma|_{[s,t]}\bigr)
=\int_s^t|\gamma'|_{\mathrm d}(r)\,\mathrm dr.
\]
若 \(X=\mathbb R^m\)，则 \(\gamma'\) 几乎处处存在，且
\[
|\gamma'|_{\mathrm d}=\|\gamma'\|
\quad\text{几乎处处},
\quad
\ell(\gamma)=\int_a^b\|\gamma'(t)\|\,\mathrm dt.
\]
\end{theorem}
\begin{proof}
度量导数定理的证明已经得到 \(|\gamma'|_{\mathrm d}=V'\) 几乎处处。由 \(V\) 的绝对连续性及子区间可加性，
\[
\ell\bigl(\gamma|_{[s,t]}\bigr)
=V(t)-V(s)=\int_s^tV'(r)\,\mathrm dr,
\]
这就是长度公式。若像在 \(\mathbb R^m\) 中，每个坐标函数都是实值绝对连续函数，因而可微的共同满测度集只需除去有限个零集。在此集合上，向量差商收敛到坐标导数组成的 \(\gamma'\)，取范数便得 \(|\gamma'|_{\mathrm d}=\|\gamma'\|\)。
\end{proof}

绝对连续性不能由“连续且可求长”替换。Cantor 函数作为 \([0,1]\) 中的一条单调连续曲线，长度为 \(1\)，度量速度却几乎处处为零。它把一部分行进过程集中在零测参数集上，导数的 Lebesgue 积分不能恢复这部分长度。

\begin{proposition}
\label{b1:prop:arc-length-parameterization}
设 \(\gamma\in\operatorname{AC}([a,b];X)\)，且 \(L=\ell(\gamma)>0\)。则存在 \(1\)-Lipschitz 曲线 \(\eta:[0,L]\to X\)，满足
\[
\gamma(t)=\eta\bigl(V(t)\bigr),
\quad
\ell\bigl(\eta|_{[u,v]}\bigr)=v-u
\quad(0\leq u<v\leq L).
\]
因而 \(|\eta'|_{\mathrm d}=1\) 几乎处处。这称为曲线的\term[huchang-canshuhua]{弧长参数化}[arc-length parameterization]。
\end{proposition}
\begin{proof}
函数 \(V\) 连续、单调不减，从 \([a,b]\) 映满 \([0,L]\)。定义 \(\eta(s)=\gamma(t)\)，其中 \(V(t)=s\)。若两个参数有相同的 \(V\) 值，则它们的像点距离不超过 \(V\) 的增量零，所以定义与选择无关。同一个位移估计给出 \(\mathrm{d}\bigl(\eta(u),\eta(v)\bigr)\leq v-u\)，故 \(\eta\) 是 \(1\)-Lipschitz 的，特别有 \(\ell(\eta)\leq L\)。

反过来，把 \(\gamma\) 的任意分割经 \(V\) 映到 \([0,L]\)，删去重复参数，距离和保持不变。因此 \(\ell(\gamma)\leq\ell(\eta)\)，遂有 \(\ell(\eta)=L\)。长度对相邻子区间可加，而每个子区间的长度都不超过其参数长度，于是 \([0,u]\)、\([u,v]\)、\([v,L]\) 三段必须分别取到这个上界。故 \(\ell\bigl(\eta|_{[u,v]}\bigr)=v-u\)。由长度公式，\(\int_u^v|\eta'|_{\mathrm d}=v-u\) 对所有子区间成立，再用 Lebesgue 微分即得 \(|\eta'|_{\mathrm d}=1\) 几乎处处。
\end{proof}

若 \(L=0\)，曲线由位移估计必为常值。对 \(L>0\)，弧长参数化与上一章的 Lipschitz 测度估计还给出
\[
\mathcal H^1\bigl(\gamma([a,b])\bigr)
=\mathcal H^1\bigl(\eta([0,L])\bigr)
\leq\mathcal H^1([0,L])=L.
\]
这里用到了\cref{b1:prop:hausdorff-lipschitz}及\cref{b1:thm:hausdorff-lebesgue}。不等式可以严格：曲线先从 \(0\) 匀速走到 \(1\)，再原路返回 \(0\)，其像为 \([0,1]\)，一维测度为 \(1\)，但长度为 \(2\)。像集测度不计重复经过，长度则把它计入；下一章的面积公式将以重数处理这种区别。

% 本章小结（不计入编号）
\begin{chaptersummary}
本章先从距离控制得到延拓和测度估计，再由一维积分表示取得沿直线的导数。
平移下的测试积分把不同方向的导数识别为同一个线性映射，
Lipschitz 条件提供的方向误差界则把有限个方向的控制提升为全部小增量的控制，
最终得到 Rademacher 定理。

\ProseParagraphBreak

近似微分允许忽略密度为零的异常点，但在局部 Lipschitz 函数上与普通微分一致。
对一般可测映射，几乎处处近似可微性可以用可数个紧集上的 Lipschitz 限制表述；
它不等于全局 Lipschitz 性，也不自动控制遗漏零测集的像。
实值绝对连续函数恰好由可积导数恢复，为后续一维 Sobolev 代表奠定基础。
在选读部分，一般绝对连续曲线展示了没有线性结构时仍可保留的部分：度量导数几乎处处存在，其积分等于行进长度。
下一章将把有限维微分矩阵的体积伸缩与映射重数结合，建立面积公式。
\end{chaptersummary}


\BookExercises

% \chapref{b1:ch:16}习题临时稿；不另设 BookExercises 标题。
% 九题均为依据本章工具组织的自拟变式，具体训练目的在各题后注明。
% 不使用\chapref{b1:ch:17}面积公式，不将未读取的教材习题冒充直接来源。

标有 \(\star\) 的习题或分项依赖\secref{b1:sec:1605}的度量曲线理论；
其余习题只使用本章主线内容。最后一题分三步重做 Rademacher 定理的关键论证。

\begin{exercise}\label{b1:ex:16-extremal-extensions}
设 \((X,\mathrm{d})\) 是度量空间，\(A\subset X\) 非空，\(f:A\to\mathbb R\) 为 \(L\)-Lipschitz 函数。令
\[
 F_-(x)=\sup_{a\in A}\bigl(f(a)-L\cdot \mathrm{d}(x,a)\bigr),
 \quad
 F_+(x)=\inf_{a\in A}\bigl(f(a)+L\cdot \mathrm{d}(x,a)\bigr).
\]
\begin{enumerate}
\item\label{b1:item:16-extension-extrema}
证明：所有 \(L\)-Lipschitz 延拓 \(F:X\to\mathbb R\) 在给定点 \(x\) 的可能取值，恰组成闭区间 \([F_-(x),F_+(x)]\)。
\item\label{b1:item:16-extension-forced-region}
取 \(X=\mathbb R\)、\(A=\{-1,1\}\)、\(f(-1)=-1\)、\(f(1)=1\)、\(L=1\)。
求 \(F_-\)、\(F_+\)，并确定所有延拓都必须取相同值的点。
\end{enumerate}
\end{exercise}
% 来源：自拟；在正文延拓存在性之外，求所有延拓的逐点取值范围及强制唯一的区域。
\begin{solution}
\secref{b1:sec:1601}的实值延拓公式给出 \(F_+\) 是一个 \(L\)-Lipschitz 延拓。
对 \(-f\) 应用同一公式再取负号，便得 \(F_-\) 也具有这一性质。
若 \(F\) 是任意这样的延拓，则对每个 \(a\in A\)，
\[
 f(a)-L\cdot \mathrm{d}(x,a)\le F(x)\le f(a)+L\cdot \mathrm{d}(x,a),
\]
取上、下确界得到 \(F_-(x)\le F(x)\le F_+(x)\)。
反过来，对 \(0\le\theta\le1\)，函数 \((1-\theta)F_-+\theta F_+\) 仍是 \(L\)-Lipschitz 延拓，其在固定点 \(x\) 的值遍历上述区间。
因此，一个点的延拓值被完全确定，当且仅当两个包络在该点相等。

在第二问中直接比较两个候选值，得到
\[
 F_-(x)=
 \begin{cases}
 x,&x\le1,\\
 2-x,&x\ge1,
 \end{cases}
 \quad
 F_+(x)=
 \begin{cases}
 -x-2,&x\le-1,\\
 x,&x\ge-1.
 \end{cases}
\]
两者恰在 \([-1,1]\) 上相等，此处所有延拓都等于 \(x\)；在区间外，每点均有至少两种可能取值。
这说明延拓在未给定数据的部分也可能被距离约束唯一确定。
\end{solution}

\begin{exercise}\label{b1:ex:16-invertible-derivative-not-local-injective}
对 \(c\ge0\)，在 \([-1,1]\) 上令 \(f_c(0)=0\)、
\(f_c(x)=x+cx^2\sin(1/x)\)（\(x\ne0\)）。
证明每个 \(f_c\) 都是 Lipschitz 函数且 \(f_c'(0)=1\)，并确定它在原点附近
具有单射邻域、具有双 Lipschitz 邻域的全部参数 \(c\)。特别处理临界参数 \(c=1\)，
不能只把正文 \(c=2\) 的反例重新计算一遍。
\end{exercise}
% 来源：自拟；以正文反例为起点，新增扰动幅度阈值及临界情形分类。
\begin{solution}
原点差商为 \(1+ch\sin(1/h)\to1\)，而非零点的导数为
\[
 f_c'(x)=1+2cx\sin(1/x)-c\cos(1/x).
\]
在区间内其绝对值不超过 \(1+3c\)，由中值定理及端点连续性得到 Lipschitz 性。
若 \(0\le c<1\)，选 \(0<\delta<1\)，使
\(2c\delta\le(1-c)/2\)。则在 \((-\delta,\delta)\) 内
\[
 (1-c)/2\le f_c'(x)\le1+3c,
\]
原点也满足同一界。中值定理给出严格正的下距离估计及上估计，故有双 Lipschitz 邻域。

若 \(c>1\)，在 \(x_n=(2\pi n)^{-1}\) 处有 \(f_c'(x_n)=1-c<0\)。
与正文例同样，连续单射必严格单调，原点的正导数排除递减，而这些负导数排除递增。

临界时取 \(s_n=1/(\pi n)\)、\(x_n=(2\pi n-s_n)^{-1}\)。
由于 \(0<s_n<1\)，有 \(1-\cos s_n\le s_n^2/2\) 及
\(\sin s_n\ge s_n-s_n^3/6\)，故
\[
 f_1'(x_n)=1-\cos s_n-\frac{2\sin s_n}{2/s_n-s_n}
 \le\frac{s_n^2}{2}
     -\frac{s_n^2(1-s_n^2/6)}{1-s_n^2/2}<0.
\]
同样不能在任何原点邻域单射。综上，两种邻域存在的参数范围都恰为 \(0\le c<1\)。
临界项的导数虽不会在 \((2\pi n)^{-1}\) 处为负，稍移开这些点后仍产生下降；
因此仅检验余弦取极值的那列点会漏掉临界情形。
\end{solution}

\begin{exercise}\label{b1:ex:16-cantor-function-not-ac}
取\cref{b1:eg:cantor-function}中 \(F_C\) 在 \([0,1]\) 上的限制。
利用正文已经证明的单调性与端点值，按区间划分的定义计算其总变差。
再利用每个第 \(n\) 级 Cantor 基本区间的 \(\rho_C\)-质量为 \(2^{-n}\)，
直接从绝对连续性的定义证明 \(F_C\) 不是绝对连续函数：写出一族两两不交的区间，
使其总长度趋于零，而对应的函数增量之和保持为 \(1\)。
\end{exercise}
% 来源：经典 Cantor 测度工具的自拟变式；不重证第14章的基础性质，只练划分变差与AC定义。
\begin{solution}
对任意划分 \(0=t_0<t_1<\cdots<t_m=1\)，单调性给出
\[
 \sum_{j=1}^m|F_C(t_j)-F_C(t_{j-1})|
 =F_C(1)-F_C(0)=1.
\]
对所有划分取上确界，得到总变差为 \(1\)。

将第 \(n\) 级的 \(2^n\) 个基本闭区间写为 \([a_j,b_j]\)。这些区间两两分离，总长度为 \((2/3)^n\)，而
\[
 F_C(b_j)-F_C(a_j)=\rho_C((a_j,b_j])=2^{-n},
\]
其中用到端点的点质量为零。因此
\[
 \sum_{j=1}^{2^n}(b_j-a_j)=(2/3)^n\longrightarrow0,
 \quad
 \sum_{j=1}^{2^n}|F_C(b_j)-F_C(a_j)|=1.
\]
取绝对连续性定义中的 \(\varepsilon=1/2\)，无论给定多小的 \(\delta>0\)，足够大的 \(n\) 都使这组区间违反定义。
这直接定位了绝对连续性失败的原因：很短的区间族仍能承载全部函数增量。
\end{solution}

\begin{exercise}\label{b1:ex:16-approximate-not-ordinary}
给定 \(\alpha>0\)、\(\beta>1\)。令 \(a_n=2^{-n}\)、\(b_n=a_n^\beta/8\)，
选 \(N\) 使 \(b_n\le a_n/16\) 对 \(n\ge N\) 成立，并定义
\[
 f_{\alpha,\beta}(x)=\sum_{n\ge N}
 a_n^\alpha\max\left\{1-\frac{|x-a_n|}{b_n},0\right\}.
\]
证明它连续且在原点近似可微，近似微分为零；分别确定它在原点普通可微、
以及在原点某邻域 Lipschitz 的参数范围。解释峰高与半宽在这两个判别中各控制什么。
\end{exercise}
% 来源：自拟；正文固定窄峰模型之后，分类峰高与峰宽两个独立指数，不重复固定参数证明。
\begin{solution}
支撑区间 \(J_n=(a_n-b_n,a_n+b_n)\) 两两分离，且只向原点聚集。
非零点附近是有限个连续峰；在 \(J_n\) 上，
\(0\le f_{\alpha,\beta}(x)\le a_n^\alpha\le(16x/15)^\alpha\)。
因为 \(\alpha>0\)，原点也连续。若 \(J_n\) 与 \((-r,r)\) 相交，则 \(a_n<2r\)，所以
\[
 m_1\left(\bigcup_{n\ge N}J_n\cap(-r,r)\right)
 \le\frac14\sum_{a_n<2r}a_n^\beta
 \le\frac{(2r)^\beta}{4(1-2^{-\beta})}.
\]
除以 \(2r\) 趋零，所有非零误差商都在这个稀疏支撑内，因而近似微分总为零。

若 \(\alpha>1\)，峰内的 \(|f_{\alpha,\beta}(x)/x|\) 被常数乘
\(x^{\alpha-1}\) 控制，峰外为零，所以普通导数存在且为零。
若 \(\alpha=1\)，峰中心差商为 \(1\)，负轴差商为零，故不可微；
若 \(0<\alpha<1\)，峰中心差商 \(a_n^{\alpha-1}\to\infty\)，同样不可微。

每个峰侧边斜率的绝对值为 \(8a_n^{\alpha-\beta}\)。
若 \(\alpha<\beta\)，这些斜率无界，排除任何原点邻域的 Lipschitz 控制。
若 \(\alpha\ge\beta\)，全部侧边斜率被一个有限数 \(M\) 控制。
同一个峰内即有 \(|f(x)-f(y)|\le M|x-y|\)。
若两点不在同一个峰内，则在它们之间有一个零值点 \(z\)；
将各点到其所在峰面向 \(z\) 的零值端点的距离记为 \(d_x,d_y\)
（峰外点取距离零），有 \(d_x+d_y\le|x-y|\)，故
\[
 |f(x)-f(y)|\le |f(x)|+|f(y)|\le M(d_x+d_y)\le M|x-y|.
\]
因此这时甚至全局 Lipschitz。普通微分的阈值是 \(\alpha>1\)，
Lipschitz 的阈值是 \(\alpha\ge\beta\)，而近似微分的稀疏性由 \(\beta>1\) 保证。
\end{solution}

\begin{optionalexercise}\label{b1:ex:16-banach-and-metric-derivatives}
在 Banach 空间 \(C([0,1])\) 中取上确界范数，并定义曲线
\[
 \gamma(t)(u)=|t-u|
 \quad(0\le t,u\le1).
\]
\begin{enumerate}
\item\label{b1:item:16-c-curve-no-derivative}
证明 \(\|\gamma(t)-\gamma(s)\|_\infty=|t-s|\)，故其度量速度在 \((0,1)\) 上处处为 \(1\)。
但是，对每个 \(t\in(0,1)\)，证明右侧差商 \((\gamma(t+h)-\gamma(t))/h\) 在 \(C([0,1])\) 中都不收敛。
\item\label{b1:item:16-l2-curve-derivative}
把同样的函数视为 \(L^2([0,1])\) 的元素，证明曲线在每个 \(t\in(0,1)\) 都可微，并写出其导数。
解释为何第一问不与绝对连续度量曲线的速度定理矛盾。
\end{enumerate}
\end{optionalexercise}
% 来源：自拟；对比同一曲线在两个范数中的微分行为，不沿用正文 L1 指示函数曲线。
\begin{solution}
反三角不等式给出 \(\bigl||t-u|-|s-u|\bigr|\le|t-s|\)，而取 \(u=0\) 即取得等号，故曲线为等距映射。
于是它是 Lipschitz 曲线，且
\[
 \frac{\|\gamma(t+h)-\gamma(t)\|_\infty}{|h|}=1.
\]

固定 \(t\in(0,1)\)。当 \(h\downarrow0\) 时，右侧差商在各个 \(u\) 点的极限为
\[
 v_t^+(u)=
 \begin{cases}
 1,&u\le t,\\
 -1,&u>t.
 \end{cases}
\]
若差商在上确界范数中收敛，其极限作为连续函数的一致极限必须连续，同时必须逐点等于 \(v_t^+\)。
这个函数在 \(t\) 不连续，矛盾。因此连右侧 Banach 空间导数也不存在。

在 \(L^2\) 中，则可忽略单点 \(u=t\)。不论 \(h\) 从哪一侧趋于零，差商在 \(u\ne t\) 处都趋于
\[
 v_t(u)=
 \begin{cases}
 1,&u<t,\\
 -1,&u>t.
 \end{cases}
\]
其在 \(u=t\) 的值可任取。差商的绝对值始终不超过 \(1\)，故与 \(v_t\) 之差的平方不超过 \(4\)。
由控制收敛定理，差商在 \(L^2\) 范数中趋于 \(v_t\)，这就是导数。
第一问的度量速度只要求距离与时间差之比有极限，不要求差商作为向量收敛；两种结论不可互换。
\end{solution}

\begin{optionalexercise}\label{b1:ex:16-retracing-multiplicity}
给定正整数 \(N\)，在节点 \(t_j=j/(2N)\)、\(0\le j\le2N\) 上设
\(\gamma_N(t_j)=0\)（\(j\) 偶数）或 \(\gamma_N(t_j)=1\)（\(j\) 奇数），并在相邻节点之间线性插值。
\begin{enumerate}
\item\label{b1:item:16-retracing-length}
求曲线 \(\gamma_N:[0,1]\to\mathbb R\) 的长度和像集的一维 Hausdorff 测度。
\item\label{b1:item:16-retracing-count}
令 \(n_N(y)=\#\{t\in[0,1]:\gamma_N(t)=y\}\)。求 \(n_N\)，并直接验证
\[
 \ell(\gamma_N)=\int_{\mathbb R}n_N(y)\,\mathrm dy.
\]
解释为何这里的计算尚不能作为任意 Lipschitz 映射的面积公式的证明。
\end{enumerate}
\end{optionalexercise}
% 来源：自拟；在正文一次往返例之外计算全部纤维重数，展示长度计入重复经过的次数。
\begin{solution}
每段的斜率绝对值为 \(2N\)，故曲线为 \(2N\)-Lipschitz，任何划分的弦长之和至多为 \(2N\)。
取包含全部节点的划分，恰有 \(2N\) 段，每段的像端点相距 \(1\)，所以长度恰为 \(2N\)。
其像始终为 \([0,1]\)，一维 Hausdorff 测度为 \(1\)。

对 \(0<y<1\)，每个线性段恰通过 \(y\) 一次，并且这些参数点均不在节点上，所以 \(n_N(y)=2N\)。
对 \(y=0\) 有 \(N+1\) 个原像，对 \(y=1\) 有 \(N\) 个原像；在 \([0,1]\) 以外原像为空。
于是 \(n_N\) 是可测函数，且
\[
 \int_{\mathbb R}n_N(y)\,\mathrm dy=2N.
\]
端点虽有不同的原像数，但它们构成 Lebesgue 零测集，不影响积分。
本题借助了有限个严格单调线性分支，重数可以直接逐段相加；一般映射没有这种现成分解，还需要建立相应的可测性和变换公式。
\end{solution}

\begin{exercise}\label{b1:ex:16-graph-differential}
设 \(f:U\to\mathbb R^k\) 为 \(L\)-Lipschitz 映射，其中 \(U\subset\mathbb R^d\) 开，并令 \(\Phi(x)=(x,f(x))\)。
\begin{enumerate}
\item\label{b1:item:16-graph-matrix}
证明 \(\Phi\) 在 \(x\) 可微当且仅当 \(f\) 在 \(x\) 可微。在这样的点，求 \(D\Phi(x)\) 及矩阵 \(D\Phi(x)^{\mathsf T}D\Phi(x)\)，并证明
\[
 1\le\det\bigl(D\Phi(x)^{\mathsf T}D\Phi(x)\bigr)
 \le(1+L^2)^d.
\]
\item\label{b1:item:16-graph-curve-length}
\(\star\)（选读）当 \(f:[a,b]\to\mathbb R^k\) 为 Lipschitz 函数时，证明其图像参数曲线的长度为
\[
 \ell(\Phi)
 =\int_a^b\sqrt{1+|f'(t)|^2}\,\mathrm dt.
\]
取 \(f(t)=|t|\)、\(-1\le t\le1\)，计算长度，并比较 \(t=0\) 处的向量导数与度量速度。
\end{enumerate}
\end{exercise}
% 来源：自拟；用微分矩阵和本章度量曲线长度公式处理图像，不使用下一章的面积公式。
\begin{solution}
若 \(f\) 在 \(x\) 可微，则将其线性近似与恒等映射合并可得
\[
 D\Phi(x)=\begin{pmatrix}I_d\\Df(x)\end{pmatrix},
 \quad
 D\Phi(x)^{\mathsf T}D\Phi(x)=I_d+Df(x)^{\mathsf T}Df(x).
\]
反过来，将 \(\Phi\) 的微分接上后 \(k\) 个坐标的投影即可得到 \(f\) 的微分。
沿直线取差商还给出 \(|Df(x)v|\le L|v|\)，于是对所有 \(v\in\mathbb R^d\)，
\[
 |v|^2\le v^{\mathsf T}\bigl(I_d+Df(x)^{\mathsf T}Df(x)\bigr)v
 \le(1+L^2)|v|^2.
\]
该实对称矩阵的全部特征值均位于 \([1,1+L^2]\)，将特征值相乘即得行列式估计。
特别是图像映射的微分总有满列秩，即使 \(Df(x)\) 自身退化也一样。

一维情形下 \(\Phi\) 是 Lipschitz 曲线，故绝对连续；在 \(f\) 可微的几乎所有点，
\(\Phi'(t)=(1,f'(t))\)，其度量速度为 \(\sqrt{1+|f'(t)|^2}\)。
由\cref{b1:thm:metric-curve-length}即得长度公式。
当 \(f(t)=|t|\) 时，除 \(0\) 外导数的绝对值为 \(1\)，故长度为 \(2\sqrt2\)。
而 \(\Phi\) 在 \(0\) 没有向量导数，但
\[
 \frac{|\Phi(h)-\Phi(0)|}{|h|}
 =\frac{\sqrt{h^2+|h|^2}}{|h|}=\sqrt2
 \quad(h\ne0),
\]
所以度量速度在这个折点仍然存在。
\end{solution}

\begin{exercise}\label{b1:ex:16-ac-composition-failure}
令 \(I_n=[2^{-n},2^{1-n}]\)，\(m_n=3\cdot2^{-n-1}\)，其中 \(n\ge1\)。
定义 \(\varphi:[0,1]\to[0,1]\)：取 \(\varphi(0)=0\)，在每个 \(I_n\) 的两个端点取值 \(0\)，在中点 \(m_n\) 取值 \(1/n^2\)，并在两侧分别线性插值。
证明 \(\varphi\) 和 \(g(u)=\sqrt u\) 都是 \([0,1]\) 上的绝对连续函数，但 \(g\circ\varphi\) 的总变差为无穷，因而不是绝对连续函数。
再说明：若将外层函数换成 Lipschitz 函数，为什么不会发生这种现象？
\end{exercise}
% 来源：自拟；用可积斜率构造 AC 内层，再由外层的非有界斜率放大短区间上的增量。
\begin{solution}
相邻区间在公共端点都取值为零，而峰高 \(1/n^2\to0\)，所以 \(\varphi\) 连续。
在各个三角峰的上升半段令 \(v=2^{n+1}/n^2\)，下降半段令 \(v=-2^{n+1}/n^2\)，在剩余可数个点任取 \(v=0\)。
这是可测函数，且
\[
 \int_0^1|v(t)|\,\mathrm dt=2\sum_{n=1}^{\infty}\frac1{n^2}<\infty.
\]
每个完整的 \(I_n\) 上 \(v\) 的积分为零。如果 \(t\in I_N\)，从 \(0\) 到 \(t\) 的积分中，所有 \(n>N\) 的完整三角峰贡献相互抵消，余下的部分恰为 \(\varphi(t)\)。
绝对可积性保证这次可数求和合法，因而
\[
 \varphi(t)=\int_0^t v(u)\,\mathrm du.
\]
由\cref{b1:thm:ac-fundamental-calculus}，\(\varphi\) 绝对连续。
同样，函数 \(u\mapsto1/(2\sqrt u)\) 在 \((0,1)\) 可积，且
\[
 \sqrt t=\int_0^t\frac{\mathrm du}{2\sqrt u},
\]
所以 \(g\) 也绝对连续。

复合函数在 \(I_n\) 两端仍为零，在中点的值却变为 \(1/n\)。
取一个包含前 \(N\) 个三角峰的端点与中点的有限划分，其增量绝对值之和至少为
\[
 2\sum_{n=1}^N\frac1n\longrightarrow\infty.
\]
故复合函数没有有限总变差，而绝对连续函数必有有限总变差。

若外层 \(G\) 为 \(L\)-Lipschitz 函数，则对任意有限组两两不交区间，
\[
 \sum_j\Bigl|G\bigl(\varphi(b_j)\bigr)-G\bigl(\varphi(a_j)\bigr)\Bigr|
 \le L\sum_j|\varphi(b_j)-\varphi(a_j)|,
\]
内层函数的绝对连续性立即传给复合函数。
本例中 \(g(u)/u=u^{-1/2}\) 在 \(u\downarrow0\) 时无界，故 \(g\) 并非 Lipschitz；正是在小峰接近零值的区域，外层映射把可求和的增量放大成了不可求和的增量。
\end{solution}

\begin{exercise}[Rademacher 定理证明的三个步骤]
\label{b1:ex:16-rademacher-interfaces}
设 \(f:\mathbb R^n\to\mathbb R\) 为 \(L\)-Lipschitz 函数，
\(V\subseteq S^{n-1}\) 是包含全部坐标方向的可数稠密集。
对 \(v\in V\)，在方向导数 \(D_vf\) 存在处令 \(g_v=D_vf\)，在其余点令 \(g_v=0\)。
\begin{enumerate}
\item\label{b1:item:16-rademacher-test}
先用一维 Lipschitz 积分表示与 Fubini 定理说明 \(D_vf\) 几乎处处存在；
再从差商的平移换元出发，证明对每个 \(\phi\in C_c^1(\mathbb R^n)\) 都有
\[
 \int_{\mathbb R^n}g_v\phi
 =-\int_{\mathbb R^n}fD_v\phi.
\]
\item\label{b1:item:16-rademacher-linearity}
记 \(g_i=g_{e_i}\)。证明
\[
 g_v=\sum_{i=1}^n v_i g_i
 \quad\text{几乎处处},
\]
并说明为什么 \(V\) 的可数性允许选出一个公共满测度集，使这些等式及全部 \(v\in V\) 的方向导数同时成立。
\item\label{b1:item:16-rademacher-net}
固定上述公共满测度集中的一点 \(x\)，令
\[
 A(h)=\sum_{i=1}^n h_i\cdot g_i(x).
\]
利用球面的有限方向网及 Lipschitz 估计，直接证明
\[
 \lim_{h\to0}\frac{|f(x+h)-f(x)-A(h)|}{|h|}=0.
\]
\end{enumerate}
\end{exercise}
% 来源：自拟；把 Rademacher 主证明的分布恒等式、共同零集与有限方向网拆成可独立复核的三步。
\begin{solution}
固定 \(v\in V\)，并令
\[
 \Delta_{v,t}f(x)=\frac{f(x+tv)-f(x)}t
 \quad(t\ne0).
\]
沿每条平行于 \(v\) 的直线，\(f\) 的限制都是一维 \(L\)-Lipschitz 函数，
故由\cref{b1:thm:one-dimensional-lipschitz-ftc}几乎处处可导。
在以 \(v\) 为一条坐标轴的正交坐标系中应用 Fubini 定理，
便知 \(D_vf\) 在 \(\mathbb R^n\) 中几乎处处存在。
平移换元给出
\[
 \int\Delta_{v,t}f(x)\cdot\phi(x)\,\mathrm dx
 =\int f(x)\frac{\phi(x-tv)-\phi(x)}t\,\mathrm dx.
\]
左端被 \(L|\phi|\) 支配，并在方向导数存在处趋于 \(g_v\phi\)。
右端只涉及 \(\operatorname{supp}\phi\) 的一个固定紧邻域，\(f\) 在其上有界；
测试函数差商又由 \(\|\nabla\phi\|_\infty\) 控制，并逐点趋于 \(-D_v\phi\)。
对两边应用控制收敛定理，就得到第一问的恒等式。

由于 \(D_v\phi=\sum_i v_i\partial_i\phi\)，第一问分别用于 \(v\) 与各坐标方向后得到
\[
 \int\left(g_v-\sum_{i=1}^n v_i g_i\right)\phi=0
 \quad\bigl(\phi\in C_c^1(\mathbb R^n)\bigr).
\]
由\cref{b1:lem:compact-tests-separate}，括号内的局部可积函数几乎处处为零。
每个 \(v\in V\) 只产生一个零测例外集；把这些集合以及各方向导数不存在的零测集作可数并，
仍得到一个零测集。因此在它的补集上，第二问中的全部关系同时成立。

现固定这样的 \(x\)。由 \(|D_vf(x)|\leq L\) 对稠密的 \(v\in V\) 成立，
线性映射 \(A\) 满足 \(\|A\|\leq L\)。给定 \(\varepsilon>0\)，取
\[
 0<\eta<\frac{\varepsilon}{2(L+\|A\|+1)},
\]
并从 \(V\) 中选出球面的有限 \(\eta\)-网 \(v_1,\ldots,v_N\)。
对这有限多个方向，\(D_{v_j}f(x)=A(v_j)\)，故存在共同的 \(\delta>0\)，使
\[
 \left|\frac{f(x+tv_j)-f(x)}t-A(v_j)\right|<\frac\varepsilon2
 \quad(1\leq j\leq N,\ 0<t<\delta).
\]
若 \(0<|h|<\delta\)，写 \(h=tu\)，其中 \(t=|h|\)、\(u\in S^{n-1}\)，并取
\(|u-v_j|<\eta\)。于是
\begin{align*}
 \frac{|f(x+h)-f(x)-A(h)|}{|h|}
 &\leq \frac{|f(x+tu)-f(x+tv_j)|}{t}\\
 &\quad+\left|\frac{f(x+tv_j)-f(x)}t-A(v_j)\right|
       +|A(u-v_j)|\\
 &\leq (L+\|A\|)\eta+\frac\varepsilon2
 <\varepsilon.
\end{align*}
这正是 \(f\) 在 \(x\) 处具有 Fr\'echet 微分 \(A\) 的定义。
\end{solution}
